% file: main.tex
\documentclass[12pt, a4paper, twoside]{article}

%\usepackage{mypackages}
%\usepackage{mycommands}
%\usepackage{myenviroments}





% PACKAGES %

%--Packages--
\usepackage[english]{babel} %if you want index titles and bibliography titles in Swedish comment this line and uncomment the following line
%\usepackage[swedish]{babel}
\usepackage[utf8]{inputenc}
\usepackage[T1]{fontenc}
%\usepackage{showkeys} %it show the label you have given to environment equation. Useful when you are writing. Remember to uncomment it before submitting
\usepackage{latexsym,amsfonts,amsmath,amsthm,bm,lmodern,titlesec,xcolor,graphicx,setspace,pdfpages,
	geometry,marginnote,caption,array,sectsty,siunitx,eurosym,colortbl,float,wrapfig,vmargin,tikz,hyperref}
\usepackage[nottoc]{tocbibind}
\usepackage{amssymb}
\usepackage{enumerate}
\usepackage{lipsum}
\usepackage{ytableau}
\usepackage{physics}
\usepackage{arydshln}
%\usepackage{subfiles}
%\usepackage{bbold}
\usepackage{tikz-cd}

%\usepackage{latexsym,amsmath,amsthm,bm,lmodern,titlesec,xcolor,graphicx,setspace,pdfpages,
	%	geometry,marginnote,caption,array,sectsty,siunitx,eurosym,colortbl,float,wrapfig,vmargin,tikz,hyperref}
%\usepackage[nottoc]{tocbibind}


\newcommand{\N}{\mathbb{N}}
\newcommand{\Z}{\mathbb{Z}}
\newcommand{\Q}{\mathbb{Q}}
\newcommand{\R}{\mathbb{R}}
\newcommand{\C}{\mathbb{C}}
\newcommand{\K}{\mathbb{K}}
\newcommand{\F}{\mathbb{F}}
\newcommand{\GL}{\text{GL}}
\newcommand{\SL}{\text{SL}}
\newcommand{\Sym}{\text{Sym}}
\newcommand{\Hom}{\text{Hom}}
\newcommand{\SO}{\text{SO}}
\newcommand{\SU}{\text{SU}}
\newcommand{\1}{\mathbf{1}}
\newcommand{\0}{\mathbf{0}}
\newcommand{\id}{\text{id}}
\newcommand{\Stab}{\text{Stab}}
\newcommand{\Cent}{\text{Cent}}
\newcommand{\Norm}{\text{Norm}}
\newcommand{\Type}{\text{Type}}
\newcommand{\D}{\text{D}}
\newcommand{\sgn}{\text{sgn}}


\newcommand{\bas}{\mathbf{\hat{e}}}
\newcommand{\ebas}{\mathbf{e}}
\newcommand{\vvec}{\mathbf{v}}
\newcommand{\uvec}{\mathbf{u}}
\newcommand{\wvec}{\mathbf{w}}
\newcommand{\ehat}{\mathbf{\hat{e}}}
\newcommand{\ihat}{\mathbf{\hat{\imath}}}
\newcommand{\jhat}{\mathbf{\hat{\jmath}}}
\newcommand{\ghat}{\mathbf{\hat{g}}}
\newcommand{\hhat}{\mathbf{\hat{h}}}
\newcommand{\vhat}{\mathbf{\hat{v}}}
\newcommand{\what}{\mathbf{\hat{w}}}
\newcommand{\uhat}{\mathbf{\hat{u}}}
\newcommand{\Span}{\text{Span}}
\newcommand{\Ss}{\scriptstyle}

\setpapersize{A4}
\setmarginsrb{30mm}{20mm}{30mm}{30mm}%left top right bottom
{8mm}{6mm}{10mm}{10mm}%headheight headsep footheight footskip


% Some useful colors
\definecolor{Rosso}{RGB}{232,61,61}
\definecolor{Viola}{RGB}{188,30,188}
\definecolor{Celeste}{RGB}{34,186,211}
\definecolor{Arancione}{RGB}{232,181,19}
\definecolor{Blu}{RGB}{0, 44, 92}
\definecolor{Azzurro}{RGB}{57,186,238}


\hypersetup{colorlinks=true,linkcolor=Blu,allbordercolors=white, urlcolor=blue, citecolor=Viola}
%\usepackage{enumitem} 
\renewcommand\labelitemi{\textbullet} 

%--Heading of  section--
\titleformat{\section}[display]
{\normalfont\bfseries\color{Blu}}
{\filleft
	\begin{tikzpicture}
		\node[
		outer sep=0pt,
		text width=2.5cm,
		minimum height=3cm,
		fill=white,
		font=\color{Blu}\fontsize{80}{0}\selectfont,
		align=center
		] (num) {\thesection};
		\node[
		rotate=90,
		anchor=south,
		font=\color{black}\large\normalfont
		] at ([xshift=-5pt]num.west) {\Huge\textsc{}};  
	\end{tikzpicture}%
}
{10pt}
{\titlerule\vskip4pt\Huge\sc\sffamily}

%--Heading of subsection--
\newcommand*\numb[1]{%
	\begin{tikzpicture}[baseline=-0.7ex]
		\node[
		outer sep=0pt,
		text width=0.6cm,
		minimum height=0.6cm,
		fill=Blu,
		font=\color{white}\fontsize{12}{20}\selectfont,
		align=center
		] (num) {\thesubsection};
	\end{tikzpicture}%
}
\titleformat{\subsection}
{\normalfont\color{Blu}\large\sc\sffamily}{\numb{\thesubsection}}{0.8em}{}

%--Heading subsubsection--
\newcommand*\stocaz[1]{%
	\begin{tikzpicture}[baseline=-0.7ex]
		\node[
		outer sep=0pt,
		text width=0.8cm,
		minimum height=0.5cm,
		fill=Blu,
		font=\color{white}\fontsize{10}{20}\selectfont,
		align=center
		] (num) {\thesubsubsection};
	\end{tikzpicture}%
}
\titleformat{\subsubsection}
{\normalfont\color{Blu}\bfseries\sc\sffamily}{\stocaz{\thesubsubsection}}{0.8em}{}

% math environments
%  we use \newtheorem to specify our theorem-like environments
% (theorem, definition, etc.) and how to display and number them.
%
% Note: The \theoremstyle declarations affect the appearance of the
% Theorems, Definitions, etc.
% Note 2: If you want swedish names you have to go and change the argument in brackets e.g. \newtheorem{thm}[defi]{Theorem} should become \newtheorem{thm}[defi]{Sats}
\theoremstyle{definition}
\newtheorem{definition}{Definition}[section]

\newtheorem*{problem}{Problem}
\newtheorem{conjecture}[definition]{Conjecture}

\theoremstyle{plain}
\newtheorem{theorem}[definition]{Theorem}
\newtheorem{proposition}[definition]{Proposition}
\newtheorem{corollary}[definition]{Corollary}
\newtheorem{lemma}[definition]{Lemma}

\theoremstyle{remark}
\newtheorem*{question}{Question}
\newtheorem{remark}[definition]{Remark}
\newtheorem{example}[definition]{Example}
\newtheorem*{notation}{Notation}
\newtheorem*{note}{Note}

\onehalfspacing
\numberwithin{equation}{section}

\usepackage{thmtools}
\renewcommand{\listoftheorems}{List of theorems and definitions}


\begin{document}
	
%	\setlength{\voffset}{0cm}\setlength{\hoffset}{0cm}\includepdf[pages={1,2,3,4}]{cover/cover}\setlength{\voffset}{-2.54cm}\setlength{\hoffset}{-2.54cm} % Saknar nummer
	
%	\clearpage{\thispagestyle{empty}}
\begin{center}
	\subsubsection*{Abstract}
	\addcontentsline{toc}{section}{Abstract}
\end{center}

	This text serves as an introduction to representation theory of finite groups, beginning with a background in group theory and linear algebra before formally defining representations, some representations are shown to be irreducible, while others are composed as a direct sum of the irreducibles (Maschke's Theorem~\ref{thm:maschkes}), culminating in the proof of Schur's Lemma~\ref{thm:schur}, which ensure the uniqueness of these compositions. Finally, character theory is introduced, simplifying representation theory by focusing on the trace of matrices, providing systematic methods to identify and decompose any given representation. Examples to demonstrate the these theoretical concepts are some cyclic and symmetric groups of small order. In particular, $\Cyc_3$, $\Cyc_4$, $\Cyc_5$, $\Sym_3$, and $\Sym_4$ are examined and completely decomposed and the results are presented in Tables~\ref{table:Cyc3},~\ref{tbl:cyc4},~\ref{table:Cyc5},~\ref{table:completecharS3}~and~\ref{table:completecharS4}.
		
\begin{center}
	\item\subsubsection*{Sammanfattning}
	\addcontentsline{toc}{section}{Sammanfattning}
	
\end{center}

	Den här texten ger en introduktion till  ändliga gruppers representationsteori. Utifrån en inledning med ett par bakgrundsämnen inom gruppteori och linjär algebra definieras representationer formellt. Vissa representationer visar sig vara oreducerbara, medan andra visar sig vara den direkta summan av oreducerbara representationer (Maschkes Sats~\ref{thm:maschkes}), vilket kulminerar i beviset av  Schurs Lemma~\ref{thm:schur} som visar att dessa sammansättningar är unika. Slutligen introduceras karaktärsteorin som förenklar arbetet med att analysera representationer genom att undersöka matrisers spår. Karaktärsteorin ger systematiska metoder för att identifiera och sönderdela varje given representation. För att demonstrera dessa teoretiska koncept ges några cykliska och symmetriska grupper som exempel, särskilt så undersöks och sönderdelas $\Cyc_3$, $\Cyc_4$, $\Cyc_5$, $\Sym_3$ och $\Sym_4$ helt och resultaten presenteras i Tabellerna~\ref{table:Cyc3},~\ref{tbl:cyc4},~\ref{table:Cyc5},~\ref{table:completecharS3}~och~\ref{table:completecharS4}.
	
%	Sammanfattningen på svenska hamnar här. %\lipsum[2][1-10]


\begin{center}
	\item\subsubsection*{Acknowledgements}
	\addcontentsline{toc}{section}{Acknowledgements}
\end{center}
	
%	\textit{NOT DONE}
	
	I would like to thank my supervisor Sofia Tirabassi for guiding me through the construction of this text. I also would like to thank Stefan Sigurdsson for proofreading the text.
	
%	Artificial intelligence (AI) services were not used in the process of planning, writing or typesetting this text. 
%	This text was typed with {\LaTeX} in TexStudio versions $\leq4.7.3$.
	%\lipsum*[3][1-10]


 % Ej skriven
	%create the table of contents
\clearpage{\thispagestyle{empty}}
{\hypersetup{linkcolor=black}
\tableofcontents
}
	
%	\clearpage{\thispagestyle{empty}}
\section{Introduction}

	Here, an introduction will be written. \lipsum[5-7] % Ej skriven
	\clearpage{\thispagestyle{empty}}
\section{Theoretical background}
	
	This is a place where I put any theorems or important definitions or concepts I build upon in my thesis. Will later be sorted into a proper section.

	\subsection{Groups}

\cite{DummitFoote}

%\begin{definition}[Center]
%	The center of a group $G$ is the set of elements of $G$ that commute and is denoted by $Z(G)$. 
%	\[
%	Z(G) = \{ g \in G \mid ga = ag \ \forall a \in G \}.
%	\]
%	Since $e$ is in $Z(G)$, it is a subgroup of $G$. If $G$ is abelian, $Z(G) = G$.	
%\end{definition}
%
%\begin{definition}[Centralizer and Normalizer]
%	The centralizer of a subset $A \subseteq G$ is the set of elements that commutes with every element of $A$.
%	The normalizer of a subset $A \subseteq G$ is the set of elements that leaves the set $A$ fixed  under conjugation. Clearly, both the centralizer and the normalizer are subgroups of $G$, and the centralizer is a subgroup of the normalizer.
%\end{definition}

\begin{definition}[Centralizer]
	The centralizer of an element $a$ of $G$ is the set of elements of $G$ that commutes with $a$, and is denoted by 
	\[
	\Cent_G(a) = \{ g \in G \mid ga = ag \}.
	\]
	It is non-empty since it contains the identity element. The centralizer of the entire group is the same as the center, $\Cent_G(G) = Z(G)$.
\end{definition}

%\begin{definition}[Stabilizer]
%	The stabilizer of an element $a$ of $G$ is the set of elements of $G$ that fixes $a$, and is denoted by 
%	\[
%	\Stab_G(a) = \{ g \in G \mid g \star a = a \}.
%	\]
%	It is non-empty since it contains the identity element.
%\end{definition}
	\subsection{The symmetric group $S_n$}

\cite[1.3]{DummitFoote}, \cite[1.1]{Sagan}

Recall that for a positive integer $n$, $S_n$ is the group of automorphisms\footnote{Set of bijections to itself} of the set $\{1, 2, \dots, n\}$.  

\begin{example}
	The six elements of $S_3$ are 
	\[
		e,(12),(13),(23),(123),(132).
	\]
	There are $3!=6$ elements: the identity element, three transpositions\footnote{2-cycles are called transpositions.} and two 3-cycles.
\end{example}


\subsubsection{Permutations}

	In $S_n$, there are $n!$ elements and they are called permutations, each defined by which elements of $\{1, 2, \dots, n\}$ it permutes.
	
	Let $(1)$ denote the identity element of any symmetric group.
	
	A permutation can be represented in some different ways, one is to compose it as a product of pairwise disjoint\footnote{Two cycles are disjoint if they have no numbers in common.} cycles, and another is to compose it as a product of transpositions.
	
	\begin{example}
		Let $\tau$ be the element of $S_6$ defined by
		\[
			\tau(1) = 2, \quad \tau(2) = 1, \quad \tau(3) = 4, \quad \tau(4) = 5, \quad \tau(5) = 3, \quad \tau(6) = 6.
		\]
		It can also be represented in \textbf{cycle notation} either as a composition of disjoint cycles
		\[
			\tau = (12)(345)(6) = (12)(345)
		\]
		consisting of a 1-cycle (usually omitted), a 2-cycle (called a \textbf{transposition}) and a 3-cycle, or as a composition of not necessarily disjoint transpositions
		\[
			\tau = (12)(34)(45).
		\]
	\end{example}

\subsubsection{Sign of a permutation}

	Let $\sigma$ be an element of $S_n$. Then the sign of $\sigma$ is defined as a function
	\begin{table}[hbt!]\centering\begin{tabular}{c c c c}
			$\sgn:$ & $S_n$    & $\rightarrow$ & $\{+1, -1\}$ \\
			&\rotatebox[origin=c]{90}{$\in$}&&\rotatebox[origin=c]{90}{$\in$} \\
			$\sgn:$ & $\sigma$ & $\mapsto$     & $(-1)^k$
	\end{tabular}\end{table}

	where $k$ is the number of transpositions required\footnote{A permutation constructed from an even number of transposition can never be written as a composition of an odd number of transpositions, and vice versa. Proof by induction on ``subcomponents'' and well-ordering principle.} to compose $\sigma$. If $k$ is even, the sign is +1 and the permutation is called even, and the reverse for an odd $k$. 
	
	The permutation $\tau$ above has sign -1, since it is composed of an odd number of transpositions.


\subsubsection{Cycle types and partitions}

	The \textbf{cycle type} of a permutation is an $n$-tuple $(k^{m_k})_{k=1}^n$ where $m_k$ is the number of $k$-cycles. The cycle type of the example above is
	\[
		\Type(\tau) = (1^1,2^1,3^1,4^0,5^0,6^0) = (1^1,2^2),
	\]
	where lastly the cycles of multiplicity 0 are omitted.
	
	Another way to describe a permutation in $S_n$ is to compare it to an integer partition of $n$. The integer partition of $n$ associated with $\sigma$ in $S_n$ is denoted by $\lambda(\sigma) = ( \lambda_1, \lambda_2, \dots, \lambda_l)$ with $l \leq n$, $\lambda_i \leq \lambda_{i+1}$ and $\sum_{i=1}^l \lambda_l = n$. The permutation $\tau$ from earlier corresponds to the partition $\lambda(\tau) = (3, 2, 1)$. 
	
	\ytableausetup{smalltableaux}
	A third equivalent way is to draw the Young diagram of the associated partition, for example the Young diagram of $\tau$ from above is 
	\begin{table}[hbt!]
		\centering
		\ydiagram{3,2,1}.
	\end{table}
	
	\begin{example}[$S_3$]
		Table \ref{table:S3} present the elements of $S_3$.	
		\begin{table}[hbt!]
			\centering
			\caption{The elements of $S_3$.}
			\begin{tabular}{r | c c c c c c}
				\label{table:S3}
				$S_3$ & e                & (12)           & (13)           & (23)           & (123)        & (132)        \\ \hline
				Type & $(1^3)$          & $(1^1,2^1)$    & $(1^1,2^1)$    & $(1^1,2^1)$    & $(3^1)$      & $(3^1)$      \\
				$\lambda$ & $(1,1,1)$        & $(2,1)$        & $(2,1)$        & $(2,1)$        & $(3)$        & $(3)$        \\
				Young & \ydiagram{1,1,1} & \ydiagram{2,1} & \ydiagram{2,1} & \ydiagram{2,1} & \ydiagram{3} & \ydiagram{3}
			\end{tabular}
		\end{table}
	\end{example}

\subsubsection{Conjugacy classes of $S_n$}

	Since the number of elements of $S_n$ is $n!$, for larger $n$ it becomes cumbersome to describe every element, however as we'll see, we can instead study the conjugacy classes of $S_n$. 
	
	Recall that two elements $g$ and $g'$ in $G$ are said to be \textbf{conjugate} if there exists an element $h$ in $G$ such that $hgh = g'$. ``Being conjugate'' in a group is an equivalence relation, so the equivalence classes (called \textbf{conjugacy classes}) partition $G$ into disjoint subsets. Denote the conjugacy class of an element $g$ in $G$ as $[g]$. If another element $g'$ is conjugate to $g$, then they share conjugacy class $[g] = [g']$ and both $g$ and $g'$ are said to be \textbf{representatives} of their common conjugacy class.
	
	The size of the conjugacy class $[g]$ can be calculated with the centralizer of $g$ in $G$, defined by 
	\[
	\Cent(g) = \left\{ h \in G \ \middle\vert \ hgh^{-1} = g \right\}
	\]
	and by the orbit-stabilizer theorem\cite[Thm 21.3]{Biggs}, the relationship between $\Cent(g)$ and the elements of $[g]$ is
	\[
	|[g]| = \frac{|G|}{|\Cent(g)|}.
	\]
	
	
	Returning to the symmetric group, two permutations $\sigma$ and $\tau$ share conjugacy class if and only if they are of the same cycle type\cite{Sagan}. Since the cycle type of a permutation was linked to an integer partition of the degree of the symmetric group, there are as many conjugacy classes of $S_n$ as there are integer partitions of $n$. Eg. there are 3 conjugacy classes in $S_3$ and 7 in $S_5$. Also, if the cycle type of a permutation $\sigma$ is $(k^{m_k})_{k=1}^n$, then the size of its centralizer is $|\Cent(\sigma)| = \prod_{k=1}^{n} k^{m_k}{m_k}!$\cite[Prop 1.1.1]{Sagan}. 
	
	\begin{example}[$S_5$]
		The conjugacy classes of $S_5$, along with their sizes and their types are presented in table \ref{table:S5}.
		\begin{table}[hbt!]
			\centering
			\caption{The conjugacy classes of $S_5$.}
			\begin{tabular}{r | c c c c c c c}
				\label{table:S5}
				$S_5$ & [e]                  & [12]               & [(12)(34)]       & [123]            & [1234]         & [(12)(345)]    & [12345]      \\ \hline
				$|\Cent(\sigma)|$ & $5!$                 & 12                 & 8                & 6                & 6              & 4              & 5            \\
				$|[\sigma]|$ & 1                    & 10                 & 15               & 20               & 20             & 30             & 24           \\
				Type & $(1^5)$              & $(1^3,2^1)$        & $(1^1,2^2)$      & $(1^2,3^1)$      & $(1^1,4^1)$    & $(2^1,3^1)$    & $(5^1)$      \\
				$\lambda $ & (1,1,1,1,1)          & (2,1,1,1)          & (2,2,1)          & (3,1,1)          & (4,1)          & (3,2)          & (5)          \\
				Young & \ydiagram{1,1,1,1,1} & \ydiagram{2,1,1,1} & \ydiagram{2,2,1} & \ydiagram{3,1,1} & \ydiagram{4,1} & \ydiagram{3,2} & \ydiagram{5}
	\end{tabular}\end{table}\end{example}

\subsection{The alternating group $A_n$}

	The subset of $S_n$ consisting of even permutations is called the alternating group of degree $n$ and is denoted by $A_n$. The order of $A_n$ is $n!/2$. The alternating group is a proper and normal subgroup of the symmetric group of same degree.

	\begin{example}
		The three elements of $A_3$ are 
		\[
		e,(123),(132).
		\]
		There are $3!/2=3$ elements: the identity element, and two 3-cycles. It is isomorphic to $\Z/3\Z$ and $Z_3$.
	\end{example}
	
	\begin{example}
		The alternating group of degree 4, $A_4$ has a normal subgroup
		\[
		V =\{e,(12)(34),(13)(24),(14)(23)\}
		\]
		This group is called the Klein-four group and it is isomorphic to $\Z/2\Z \times \Z/2\Z$.
	\end{example}
	
\subsection{The dihedral group $\D_{2n}$}

	The set of symmetries of a regular $n$-gon centered on the origin in the plane is called the dihedral group of order $2n$ and is denoted by $\D_{2n}$. It can be constructed via two generators $r$ and $s$ where $r$ is the rotation of the $n$-gon by $2\pi/n$ around the origin and $s$ is the reflection along the axis between the origin and an arbitrary ``first'' vertex, and the relations between them. Clearly, the order of $r$ is $n$ and $s$ is 2, so
	\[
		\D_{2n} = \left\langle r,s \ \middle\vert \ r^n = s^2 = e, sr = rs^{-1} \right\rangle.
	\] 
	
	\begin{example}
		The symmetries of the triangle is
		\[
		D_6 = \{ e, r, r^2, s, sr, sr^2 \},
		\]
		hence it is isomorphic to $S_3$.
	\end{example}
	
	
	\subsection{Group action}

	Right now, the definitions of group action of a group on a set and on itself by ``left'' multiplication is presented in later section on permutation and regular representations.

\begin{definition}[Group action]
	Group action...
\end{definition}
	\subsection{Vector spaces over a field}

In part based on \cite{Jeevanjee}.

\begin{remark}[Modules]
	If $V$ is a vector space over a field $\K$, then $V$ is the same as a $\K$-module.
\end{remark}

\begin{definition}[Dual space]
	For any vector space $V$ over the field $\K$ there is a corresponding dual vector space denoted $V^*$. If the elements of $V$ are $\{\mathbf{v}_i\}_{i=1}^n$ and a basis is $\{\mathbf{e}_i\}_{i=1}^n$, then the dual space $V^*$ has elements $\{\mathbf{v}^i\}_{i=1}^n$ and a basis is $\{\mathbf{e}^i\}_{i=1}^n$. An element $v^i$ of $V^*$ can also be interpreted as a $(1,0)$-tensor
	\[
	v^i: V \longrightarrow \K,
	\]
	in this sense, $V^*$ is the set of linear maps from $V$ to $\K$ denoted by $V^* = \Hom(V,\K)$.
	The spaces $V$ and $V^*$ carry a ``natural pairing'' that is a bilinear mapping
	\[
	\bra{\mathbf{v}^j}\ket{\mathbf{v}_i}: V \times V^* \longrightarrow \K
	\]
	that for basis vectors is defined as
	\[
	\bra{\mathbf{e}^j}\ket{\mathbf{e}_i} = {\delta_i}^j.
	\]
\end{definition}

\begin{definition}[Orthogonal subspace]
	\cite[12.7]{holst} If $W$ is a subspace in $V$, then its complementary subspace in $V$ is defined by 
	\[
	W^\perp = \left\lbrace \vvec \in V \ \middle\vert \ \wvec^\dagger \vvec = 0, \text{ for any } \wvec \in W \right\rbrace.
	\]
\end{definition}

\begin{lemma}[Complementary subspaces]\label{lemma:compsubspace}\cite[Thm 12.16]{holst}
	Let $W$ be a subspace of a space $V$. Let $\dim V = n$ and $\dim W = k$, then:
	\noindent
	\begin{enumerate}[i)]
		\setlength\itemsep{-0.5em}
		\item $\dim W^\perp = n - k$,
		\item $V = W \oplus W^\perp$,
		\item If $\{\ebas_i\}_{i=1}^{n}$ is the basis in $V$, then $\{\ebas_i\}_{i=1}^{k}$ and $\{\ebas_i\}_{i=k+1}^{n}$ are the bases for $W$ respectively $W^\perp$.
	\end{enumerate}
\end{lemma}

\subsubsection{Tensor operations on vector spaces}

	\cite{Jeevanjee}, but some from \cite{holst}
	
	In this section, multilinearity will be presented. Two tensors, operators, matrices etc. are said to be \textbf{bilinear} if when fixing either one, then the other is linear, and vice versa. This is extended to an arbitrary number of pairwise bilinear tensors, then they are called \textbf{multilinear}.
	
	In the following few definitions, let $U$ and $V$ be two vector spaces over the same field $\K$ with respective dimensions $m$, $n$ and bases $\{\uhat\}$, $\{\vhat\}$.
	
	\begin{definition}[Direct sum of vector spaces]
		The direct sum of $U$ and $V$, denoted $U \oplus V$ is a vector space over $\K$ with a basis constructed by all $\uhat \oplus \vhat$. The dimension of $U \oplus V$ is $m+n$. Let $F \in \GL(U)$ and $G \in \GL(V)$. Then $F \oplus G \in \GL(U \oplus V)$ is a block matrix
		\begin{align}
			\begin{pmatrix}
				F & \0 \\
				\0 & G
			\end{pmatrix},
		\end{align}
		where $\0$ are zero matrices ``fitting in the corners'', acting on a vector $\uvec \oplus \vvec$ of $U \oplus V$ like:
		\begin{align}
			\left(
			\begin{array}{ccc;{2pt/2pt}ccc}
				f_{11}                    & \cdots & f_{1m} & 0      & \cdots & 0      \\
				\vdots                    & \ddots & \vdots & \vdots & \ddots & \vdots \\
				f_{m1}                    & \cdots & f_{mm} & 0      & \cdots & 0      \\
				\hdashline[2pt/2pt]
				0 & \cdots & 0      & g_{11} & \cdots & g_{1n} \\
				\vdots                    & \ddots & \vdots & \vdots & \ddots & \vdots \\
				0                         & \cdots & 0      & g_{n1} & \cdots & g_{nn}
			\end{array}
			\right)\cdot
			\begin{pmatrix}
				u_1                         \\
				\vdots                          \\
				u_{m}                       \\
				\hdashline[2pt/2pt]
				v_1 \\
				\vdots                          \\
				v_{n}
			\end{pmatrix}.
		\end{align}
	\end{definition}
	
	\begin{definition}[Tensor product of vector spaces]\label{def:tensorproduct}
		The tensor product of $U$ and $W$, denoted $U \otimes V$ is a vector space over $\K$ with a basis constructed by all $\uhat \otimes \vhat$. The dimension of $U \otimes V$ is $mn$. Let $F \in \GL(U)$ and $G \in \GL(V)$. Then $F \otimes G \in \GL(U \otimes V)$ is a block matrix
		\begin{align}
			(f_{ij}G)_{m \times n}=
			\left(
			\begin{array}{c c c}
				f_{11}G & \cdots & f_{1m}G \\
				\vdots & \ddots & \vdots \\
				f_{m1}G & \cdots & f_{mm}G
			\end{array}
			\right)			
		\end{align}
		where $f_{ij}G$ is understood to be the scalar product of $f_{ij}$ with $G$.
	\end{definition}
	
	\begin{definition}[Tensor power]
		If $V$ is a vector space over $\K$, then for a positive integer $n$, the $n$th tensor power of $V$, denoted by $V^{\otimes n}$ is also a vector space, by induction on definition \ref{def:tensorproduct}.
		\begin{notation}
			The zeroth power of a vector space is defined as its ground field, that is $V^{\otimes 0} = \K$, and the first power is the vector space itself, that is $V^{\otimes 1} = V$.
		\end{notation}    
	\end{definition}
%	\clearpage{\thispagestyle{empty}}
\section{Groups}

This section on groups is mostly based on the treatment on groups of \cite{DummitFoote}.

\begin{definition}[Group]
	A group is a set $G$ together with a binary operation $\star$ fulfilling:
	\begin{enumerate}[G1.]
		\setlength\itemsep{0em}
		\item Closure under the operation.%The set is closed under the operation, that is $gh$ is also an element.
		\item Associativity under the operation.%The operation is associative, that is $g(hi) = ghi = (gh)i$.
		\item Existence of a unique identity element.%There exists a unique \textbf{identity element} $e$ in $G$ that leaves every element fixed, that is $ge = g = eg$.
		\item Existence of unique inverses.%For every element $g$ there exists a unique \textbf{inverse element} $g^-1$ such that $gg^{-1} = e = g^{-1}g$.
		\textit{\item\label{axiom:G4} (If the group is commutative, it is called abelian.)}%(If every element of $G$ commutes, that is $gh = hg$, the group is said to be commutative, or \textbf{abelian}.)}
\end{enumerate}
\end{definition}

Sometimes the group is denoted by $(G,\star)$ but usually the operation is dropped and the set $G$ alone represents the group, especially if the operation is inferred from context. The last axiom G5 is not generally true for all groups. 

\begin{example}[Additive groups]
The number sets $\Z$, $\Q$, $\R$ and $\C$ are all abelian group under ordinary addition. %Their identity element is $0$ and for a number $n$, its inverse is $-n$. 
A vector space $V$ over a field $\K$ is an abelian group under vector addition. %The identity element is the zero vector $\vec{0}$ and the inverse of a vector $v \in V$ is $-v$.
For an integer $n>0$, the set of residue classes of modulo $n$, denoted $\Z/n\Z$, is an group under addition modulo $n$.
\end{example}

\begin{example}[Multiplicative groups]
The sets $\Q$, $\R$ and $\C$ without $0$ (denoted $\Q^{\times}$, $\R^{\times}$ and $\C^{\times}$) are  groups under ordinary multiplication. %The identity element is $1$ and the inverse of the number $n$ is $1/n = n^{-1}$.
The general linear group (for an integer $n>0$, the set of invertible\footnote{Non-zero determinant} $n \times n$-matrices with entries from the field $\K$), denoted $\GL_n(K)$ is a group under matrix multiplication. 
The set $(\Z/n\Z)^{\times}$ is a group under multiplication modulo $n$.
\end{example}

\subsection{Subgroups}

\begin{definition}[Subgroup]
A non-empty subset $H \subseteq G$ is called a subgroup of $G$ if it is closed under inverses. Then it is denoted by $H \leq G$.
\end{definition}

Every group $G$ has these subgroups:

\begin{example}
The group $G$ is a subgroup of itself. However, a proper subgroup $H \neq G$ is denoted by $H < G$. The singleton $\{e\}$ is called the trivial subgroup.
\end{example}

The examples from the previous section has the following subgroups.

\begin{example}
The ordinary number sets are related by $\C \subset \R \subset \Q \subset \Z$. The integers has subgroups $n\Z$ for an integer $n>0$. The special linear group with entries from a field $\K$, denoted by $\SL_n(\K)$ is a subgroup of $\GL_n(\K)$.
\end{example}

\begin{example}[Center]
The center of a group $G$ is the set of elements of $G$ that commute and is denoted by $Z(G)$. Since $e$ is in $Z(G)$, it is a subgroup of $G$. If $G$ is abelian, $Z(G) = G$.	
\end{example}

\begin{example}[Centralizer and Normalizer]
The centralizer of a subset $A \subseteq G$ is the set of elements that commutes with every element of $A$.
The normalizer of a subset $A \subseteq G$ is the set of elements that leaves the set $A$ fixed  under conjugation. Clearly, both the centralizer and the normalizer are subgroups of $G$, and the centralizer is a subgroup of the normalizer.
\end{example}

\begin{definition}[Normal subgroup]
A subgroup $N \leq G$ is called normal if it is invariant under conjugation by $G$. It is denoted by $N \trianglelefteq G$. A proper normal subgroup is denoted by $N \triangleleft G$. 
\end{definition}

\begin{example}
Any group is a normal subgroup to itself. In an abelian group, every subgroup is normal. In $GL_n(\K)$, $\SL_n(\K)$ is normal.
\end{example}

\begin{definition}[Simple group]
A group is called simple if there are no non-trivial and proper normal subgroups.
\end{definition}

\begin{example}
The integers modulo $n$ is a simple group if $n$ is a prime. For example, $\Z/7\Z$ has only itself and $\{e\}$ as normal subgroups, and as a non-example the integers have the even integers as a non-trivial and proper normal subgroup.
\end{example}

\subsection{The symmetric group $\mathfrak{S}_n$}

The set of all permutations of the numbers $\{1,2,3,\dots,n\}$ is called the symmetric group of degree $n$ and is denoted by\footnote{The strange symbol is a ``fraktur'' $S$.} $\mathfrak{S}_n$. Equivalently, it is the set of bijections from $\{1,2,3,\dots,n\}$ to itself. The order of $\mathfrak{S}_n$ is $n!$. 

\begin{example}
The elements of $\mathfrak{S}_3$ are 
\[
e,(12),(13),(23),(123),(132).
\]
There are $3!=6$ elements: the identity element, three transpositions\footnote{2-cycles are called transpositions.} and two 3-cycles.
\end{example}


\subsection{The alternating group $\mathfrak{A}_n$}

The subset of $\mathfrak{S}_n$ consisting of even\footnote{A permutation is considered even if it can be decomposed to an even number of 2-cycles, or equivalently it is of odd length.} permutations is called the alternating group of degree $n$ and is denoted by $\mathfrak{A}_n$. The order of $\mathfrak{A}_n$ is $n!/2$. The alternating group is a normal subgroup of the symmetric group of same degree.

\begin{example}
The elements of $\mathfrak{A}_3$ are 
\[
e,(123),(132).
\]
There are $3!/2=3$ elements: the identity element, and two 3-cycles. It is isomorphic to $\Z/3\Z$.
\end{example}

\begin{example}
The alternating group of degree 4, $\mathfrak{A}_4$ has a normal subgroup
\[
V =\{e,(12)(34),(13)(24),(14)(23)\}
\]
This group is called the Klein-four group and it is isomorphic to $\Z_2^2$.
\end{example}

\subsection{Dihedral groups $\D_{2n}$}

The set of symmetries of a regular $n$-gon centered on the origin in the plane is called the dihedral group of order $2n$ and is denoted by $\D_{2n}$. It can be constructed via two generators $r$ and $s$ where $r$ is the rotation of the $n$-gon by $2\pi/n$ around the origin and $s$ is the reflection along the axis between the origin and an arbitrary ``first'' vertex, and the relations between them. Clearly, the order of $r$ is $n$ and $s$ is 2, so
\[
D_{2n} = \left\langle r,s \ \middle\vert \ r^n = s^2 = e, sr = rs^{-1} \right\rangle.
\] 

\begin{example}
The symmetries of the triangel is
\[
D_6 = \{ e, (123), (132), (12), (23), (13) \},
\]
hence it is isomorphic to $\mathfrak{S}_3$.
\end{example}

\subsection{Quotient groups}

The normal subgroups subgroups discussed above are important, because they permit the existence of quotient groups.

\begin{example}
The integers $\Z$ has a normal\footnote{Any subgroup of $\Z$ is normal since $\Z$ is abelian.} subgroup $n\Z$. The quotient group $\Z/n\Z$ are the integers modulo $n$.
\end{example}

\subsubsection{The first isomorphism theorem}

Something something.

\subsection{Group action}

Somethnig something.

\subsubsection{Action of a group on itself by multiplication}

Something something.

\subsubsection{Action of a group on itself by conjugation}

Something something.

\subsubsection{Conjugacy classes}

Something something.

 % RENSAD
%	\clearpage{\thispagestyle{empty}}
\section{Some introductory background topics}

\subsection{Group theory}



%This section on groups is mostly based on the treatment of \cite{DummitFoote}.
%
%\subsubsection{Introduction to Groups}
%
%\begin{definition}[Group]
%	A group is a set $G$ together with a binary operation\footnote{Examples of binary operators are ordinary addition or multiplication, addition $\mod n$ for some integer $n$, vector addition, composition of functions and matrix addition and multiplication.} $\star: G \times G \rightarrow G$ fulfilling:
%	\begin{enumerate}[i)]
%		\item Associativity over the binary operation. For any $a,b,c \in G$, $a \star (b \star c) = (a \star b) \star c = a \star b \star c$.
%		\item Existence of a unique identity element in $G$. There is an element $e \in G$ such that for any $a \in G$, $a \star e = a = e \star a$.
%		\item Existence of unique inverses for every element of $G$. For every $a \in G$, there is an element denoted $a^{-1} \in G$ such that $a \star a^{-1} = e = a^{-1} \star a$.
%	\end{enumerate}
%\end{definition}
%
%Usually the $\star$ symbol is dropped when it is clear from context what operation we are using. In an additive group, the operation is addition. Usually something that looks like zero\footnote{For example 0, the zero vector, the zero matrix...} is the additive identity element and an inverse of an element $a$ is usually denoted by $-a$.
%
%\begin{example}[Some additive groups]
%	Some examples of additive groups are $\Z$, $\Q$, $\R$ and $\C$; any vector space under vector addition and the integers mod $n$, denoted $\Z / n \Z$.
%\end{example}
%
%In a multiplicative group, the operation is usually some sort of multiplication or composition. Usually something that looks like one\footnote{For example 1, the unit matrix $\1_n = (\delta_{ij})_{n \times n}$, the identity permutation $\id = e$...} is the multiplicative identity element.
%
%\begin{example}[Some multiplicative groups]	
%	Some examples of multiplicative groups are $\Z^\times$, $\Q^\times$, $\R^\times$ and $\C^\times$ (these are the subsets of their additive counterpart where zero has been removed since it lacks a multiplicative inverse.); $(\Z / n \Z)^\times$; The set of all invertible square $n \times n$ matrices denoted by $\GL(\C)$; Difference between GL(V) and GL(C)?
%\end{example}
%
%In general, any two elements of a group do not commute, that is for $a,b \in G$, $ab \neq ba$. However the special case when every element is commutative is important.
%
%\begin{definition}[Abelian group]
%	A group $G$ in which every element commutes with every other element (that is $ab = ba$ for every $a,b \in G$) is called a commutative group, or an Abelian group.
%\end{definition}
%
%\subsubsection{Subgroups}
%
%\begin{definition}[Subgroup]
%	A non-empty subset $H$ of $G$ is called a subgroup of $G$ if it fulfills:
%	\begin{enumerate}[i)]
%		\item Closure under the operation of $G$. If $h_1, h_2 \in H$, then $h_1 h_2, h_2 h_1 \in H$.
%		\item Closure under inverses. If $h \in H$, then there is an element $h^{-1} \in H$.
%	\end{enumerate}
%	If $H$ is a subgroup $G$, we denote the relation by $H \leq G$. The trivial subgroup is the subset of $G$ only consisting of the identity element of $G$, that is $\{e\}$. A proper subgroup $H$ of $G$ is a proper subset of $G$ and is denoted by $H < G$.
%\end{definition}

%\begin{definition}[Centralizer]
%	The centralizer of an element $a$ of $G$ is the set of elements of $G$ that commutes with $a$, and is denoted by 
%	\[
%	\Cent_G(a) = \{ g \in G \mid ga = ag \}.
%	\]
%	It is non-empty since it contains the identity element.
%\end{definition}
%
%\begin{definition}[Stabilizer]
%	The stabilizer of an element $a$ of $G$ is the set of elements of $G$ that fixes $a$, and is denoted by 
%	\[
%	\Stab_G(a) = \{ g \in G \mid g \star a = a \}.
%	\]
%	It is non-empty since it contains the identity element.
%\end{definition}

\subsubsection{Some special and interesting groups}

%\begin{definition}[The dihedral group of order $2n$]
%	For every integer $n \geq 3$, the set of all symmetries of a regular planar $n$-gon centered on the origin is called the dihedral group of order $2n$ and is denoted by $\D_{2n}$. The vertices of the regular $n$-gon are labeled clockwise from 1 to $n$. The dihedral group is generated by the two kinds of operations permitted on this rigid $n$-gon: rotation and reflexion. We define relection $s$ about the symmetry line through vertex 1 and the origin. Rotation $r$ is defined by either moving every vertex $i$ to the vertex $i+1 \mod n$, or one step clock-wise or equivalently by rotating the $n$-gon $2 \pi /n$ radians. We see that the order\footnote{The order of an element $g \in G$ is the smallest possible positive integer $m$ such that $g^m = e$.} of $s$ is 2 and of $r$ is $n$.
%	
%	The dihedral group can be constructed using $s$ and $r$ as generators along with some information on how $s$ and $r$ relate:
%	$$ \D_{2n} = \left\langle r, s \mid r^n = s^2 = 1, rs = sr^{-1} \right\rangle.$$
%\end{definition}
%
%\begin{definition}[The symmetric group of degree $n$]
%	The set of bijections from the set $\{1, 2, \dots n\}$ to itself if denoted $S_n$ and is called the symmetric group of degree $n$. Clearly $S_n$ has order $n!$. The associated operation is cycle decomposition $\circ$. For two permutations $\sigma, \tau \in S_n$, their decomposition is $\sigma \circ \tau = \sigma \tau$.
%\end{definition}
%
%The elements of $S_n$ can be represented in some different ways, one is to decompose every permutation $\sigma \in S_n$ to a product of pairwise disjoint\footnote{Two cycles are disjoint if they have no numbers in common.} cycles. The congruence class an element of $S_n$ is determined by the shape of its cycle decomposition, for example if an element consist of a 2-cycle and a 3-cycle, it is in the same conjugacy class of every other such cycle. A representative of a congruence can be chosen arbitrarily, and is surrounded by square brackets, for example is $\left[ e \right]$ is the conjugacy class consisting of only $e$.
%
%
%\begin{example}[$S_3$ and $S_4$]
%	The elements of $S_3$ and $S_4$ are presented in the following tables.
%	
%	%			\begin{center}
%		%				\begin{tabular}{r c l }
%			%					                   & $S_3$ &                      \\
%			%					  $\left[e\right]$ & $=$   & $\{e\}$              \\
%			%					 $\left[12\right]$ & $=$   & $\{(12),(13),(23)\}$ \\
%			%					$\left[123\right]$ & $=$   & $\{(123),(132)\}$
%			%				\end{tabular}
%		%				\\
%		%				\begin{tabular}{r c l }
%			%					                        & $S_4$ &                                                       \\
%			%					       $\left[e\right]$ & $=$   & $\{e\}$                                               \\
%			%					      $\left[12\right]$ & $=$   & $\{(12),(13),(14),(23),(24),(34)\}$                   \\
			%					     $\left[123\right]$ & $=$   & $\{(123),(124),(132),(134),(142),(143),(234),(243)\}$ \\
			%					    $\left[1234\right]$ & $=$   & $\{(1234),(1243),(1324),(1342),(1423),(1432)\}$       \\
			%					$\left[(12)(34)\right]$ & $=$   & $\{(12)(34),(13)(24),(14)(23)\}$
			%				\end{tabular}
		%			\end{center}
%\end{example}

%A subgroup of $S_n$ is the subset of even\footnote{A permutation is called even(odd) if its order is even(odd) permutations.} permutations. This is called the alternating group. For the examples above,

%\begin{example}[$A_3$ and $A_4$]
%	The elements of $S_3$ and $S_4$ are presented in the following tables.
%	
	%			\begin{center}
		%				\begin{tabular}{r c l }
			%					                   & $A_3$ &                   \\
			%					  $\left[e\right]$ & $=$   & $\{e\}$           \\
			%					$\left[123\right]$ & $=$   & $\{(123),(132)\}$
			%				\end{tabular}
		%				\\
		%				\begin{tabular}{r c l }
			%					                        & $A_4$ &                                                       \\
			%					       $\left[e\right]$ & $=$   & $\{e\}$                                               \\
			%					     $\left[123\right]$ & $=$   & $\{(123),(124),(132),(134),(142),(143),(234),(243)\}$ \\
			%					$\left[(12)(34)\right]$ & $=$   & $\{(12)(34),(13)(24),(14)(23)\}$
			%				\end{tabular}
		%			\end{center}			
%\end{example}

%Clearly, $A_n$ is a subgroup of $S_n$. $S_4/S_3 \cong A_3$???

%\subsubsection{Quotient groups}
%
%Quotient group???
%
%\subsubsection{Group action}
%
%Group action
%Groups acting on themselves - conjugation - class equation
%
%Automorphisms
%
%\subsection{Rings and fields}
%
%\subsection{Vector spaces}
 % RENSAD
%	\clearpage{\thispagestyle{empty}}
\section{Some introductory topics}
\subsection{Vector spaces over a field}

%	In part based on \cite{Jeevanjee}.
%
%	\begin{definition}[Vector space over a field, dimension, basis]
%		A vector space over a field is an abelian group of vectors $V$ together with a field of scalars $\K$, equipped with vector addition and scalar multiplication. 
%		A subset of $V$ that is component-wise linearly independent and spans the vector space is called a basis and is denoted by $\left\lbrace\mathbf{e}_i\right\rbrace_{i=1}^n = \left\lbrace \mathbf{e}_1, \mathbf{e}_2, \dots, \mathbf{e}_n \right\rbrace$. There may be multiple bases in a vector space. 
%		The number of basis vectors needed to span the space is called the dimension of the space and is denoted by $n = \dim V$. 
%		The natural choice of basis for $\R^n$ is the orthonormal basis $\left\lbrace \delta_{ij} \right\rbrace_{i=1}^n$.
%	\end{definition}
%	
%	\begin{example}
%		For a positive integer $n$, $\K^n$ are vector spaces over any field $\F$. The group of invertible $n \times n$-matrices $\GL_n(\K)$ is not a vector space since $\mathbf{0}$ is not in $\GL_n(\K)$, however the set of all $n \times n$-matrices $\mathcal{M}_n(\K)$ is a vector space (it is an additive group).
%	\end{example}
	
%	\begin{remark}[Modules]
%		If $V$ is a vector space over a field $\K$, then $V$ is the same as a $\K$-module.
%	\end{remark}
	
%	\begin{definition}[Dual space]
%		For any vector space $V$ over the field $\K$ there is a corresponding dual vector space denoted $V^*$. If the elements of $V$ are $\{\mathbf{v}_i\}_{i=1}^n$ and a basis is $\{\mathbf{e}_i\}_{i=1}^n$, then the dual space $V^*$ has elements $\{\mathbf{v}^i\}_{i=1}^n$ and a basis is $\{\mathbf{e}^i\}_{i=1}^n$. An element $v^i$ of $V^*$ can also be interpreted as a $(1,0)$-tensor
%		\[
%			v^i: V \longrightarrow \K,
%		\]
%		in this sense, $V^*$ is the set of linear maps from $V$ to $\K$ denoted by $V^* = \Hom(V,\K)$.
%		The spaces $V$ and $V^*$ carry a ``natural pairing'' that is a bilinear mapping
%		\[
%			\bra{\mathbf{v}^j}\ket{\mathbf{v}_i}: V \times V^* \longrightarrow \K
%		\]
%		that for basis vectors is defined as
%		\[
%			\bra{\mathbf{e}^j}\ket{\mathbf{e}_i} = {\delta_i}^j.
%		\]
%	\end{definition} % RENSAD
%	%\clearpage{\thispagestyle{empty}}
\subsection{Tensor algebra}

	
%	
%	\begin{definition}[Tensor]
%		A $(r,s)$-tensor $T$ is a multi-linear function
%		\[
%			T: \underset{r \text{ times}}{\underbrace{V \times \dots \times V}} \times \underset{s \text{ times}}{\underbrace{V^* \times \dots \times V^*}} \longrightarrow \C.
%		\]
%		Or in other words, it takes $r$ vectors from a vector space and $s$ vectors from the corresponding dual vector space $V^*$ and maps it to a number in $\C$.
%	\end{definition}
%	
%	\begin{example}[Tensors]
%		\begin{itemize}
%			\itemsep0em 
%			\item A $(0,0)$-tensor is per definition a scalar.
%			\item A $(1,0)$-tensor is a vector.
%			\item A $(0,1)$-tensor is a dual vector.
%			\item An example of a $(3,0)$-tensor is the Levi-Civita tensor defined by 
%			\[
%				\epsilon(\mathbf{u}, \mathbf{v}, \mathbf{w}) = (\mathbf{u} \cdot \mathbf{v}) \times \mathbf{w},
%			\]
%			which could be thought of as a measurement of the volume of the parallelepiped spanned by the vectors $\mathbf{u},\mathbf{v}$ and $\mathbf{w}$. When applied on any basis vectors $\mathbf{e}_i, \mathbf{e}_j$ and $\mathbf{e}_k$ of a vector space, we find the Levi-Civita symbol
%			\[
%				\epsilon(\mathbf{e}_i, \mathbf{e}_j, \mathbf{e}_k) := \epsilon_{ijk} = \begin{cases}
%					1, \text{ if $i,j,k$ cyclical} \\
%					-1, \text{ if $i,j,k$ anti-cyclical} \\
%					0, \text{ if $i,j,k$ non-distinct.}
%					\end{cases}.
%			\]
%		\end{itemize}
%	\end{example}
	
%	\begin{definition}[Direct sum of vector spaces]
%		content...
%	\end{definition}
%	
%	\begin{definition}[Tensor product of vector spaces]\label{def:tensorproduct}
%		Let $V$ and $W$ be two vector spaces over the field $\K$. Their \textbf{tensor product} is also a vector space over $\K$ and is denoted by $V \otimes W$. Its elements will look like $\mathbf{v} \otimes \mathbf{w}$ if $\mathbf{v}$ is from $V$ and $\mathbf{w}$ is from $W$. 
%	\end{definition}
%	
%	\begin{definition}[Linear combination of vector spaces]
%		content...
%	\end{definition}
%	
%	\begin{definition}[Tensor power]
%		If $V$ is a vector space over $\K$, then for a positive but not necessarily non-zero integer $n$, the $n$th tensor power of $V$, denoted by $V^{\otimes n}$ is also a vector space, by induction on definition \ref{def:tensorproduct}.
%		\begin{notation}
%			The zeroth power of a vector space is defined as its ground field, that is $V^{\otimes 0} = \K$, and the first power is the vector space itself, that is $V^{\otimes 1} = V$.
%		\end{notation}
%	\end{definition}
	
%	\begin{example}[$\R^3 \otimes \R^3 = {\R^3}^{\otimes 2}$]
%		Choosing a basis $\{\ebas_1, \ebas_2, \ebas_3\}$ for $\R^3$, a basis of $\R^3 \otimes \R^3$ is 
%		\begin{table}[hbt!]\centering\begin{tabular}{r c c c l}
%			$\mathcal{B} = \{\{ \ebas_i \otimes \ebas_j \}_{i=1}^3\}_{j=1}^3 = \{$ & $\ebas_1 \otimes \ebas_1$, & $\ebas_1 \otimes \ebas_2$, & $\ebas_1 \otimes \ebas_3$, &  \\
%			     & $\ebas_2 \otimes \ebas_1$, & $\ebas_2 \otimes \ebas_2$, & $\ebas_2 \otimes \ebas_3$, &  \\
%			     & $\ebas_3 \otimes \ebas_1$, & $\ebas_3 \otimes \ebas_2$, & $\ebas_3 \otimes \ebas_3$ & $\}$.
%		\end{tabular}\end{table}
%	\end{example} % RENSAD
	
	\clearpage{\thispagestyle{empty}}
\section{Representation theory}

%	\textit{Detta är syntesen av tre tidigare versioner. }
	
	The layout of this section is based on \cite{Serre} and \cite{FultonHarris}. From this point on, only finite groups and finite-dimensional vector spaces are considered. 
	
	Let $G$ be a group written multiplicatively and let its order (or size or cardinality) be denoted by $|G|$. 
	
	Let $V$ be a vector space over a field $\K$. Its group of automorphisms is denoted by $\GL(V)$. If $V$ has a dimension $\dim V = n$ and we provide a basis, usually denoted $\{\bas_i\}_{i=1}^n$, then $\GL(V)$ is identified with the set of linear transformations of $V$, which is the set of invertible square matrices of size $n \times n$, denoted $\GL_n(\K)$\cite[18.1]{DummitFoote}. The main purpose of this text is to study linear representations of finite groups, a process explained as:
	
	\textbf{\emph{A representation of a group in a vector space is a homomorphism that takes an element of the group into the group of automorphisms of the vector space.}}
	
	Or more directly in the following definition.
	
	\begin{definition}[Representation]
		Given a finite group $G$, a representation of $G$ in $V$ is a homomorphism
		\begin{align}
			\rho: G \rightarrow \GL(V).
		\end{align}
	\end{definition}
	For every $g$ in $G$, there is a transformation $\rho(g)$ in $\GL(V)$, or again $\rho(g): V \rightarrow V$ for some $g$ in $G$. In other words, $\rho(g)$ acts acts as a linear map on a vector $v$ in $V$ and maps it to the vector $\rho(g)(v)$ in $V$.
	
	Recall that if $\rho$ is a homomorphism, then for any two elements $g$ and $h$ in $G$,
	\begin{align}\label{eq:homomorphism}
		\rho(gh) = \rho(g)\rho(h),
	\end{align}
	and these two consequences follows:	
	\begin{proposition}
		For the identity element $e$ of $G$ and an arbitrary element $g$ with inverse $g^{-1}$ in $G$,
		\begin{align}
			\rho(e) = \id
		\end{align}
		and
		\begin{align}
			\rho(g)^{-1} = \rho(g^{-1}),
		\end{align}
		where $\id$ is the identity transformation and $\rho(g)^{-1}$ is the inverse representation of $g$.
	\end{proposition}
	\begin{proof}
		Take $g$ as an arbitrary element of $G$. The first identity follows from taking $h=e$ in eq.\ref{eq:homomorphism}:
		\begin{align}
			\rho(eg) = \rho(e)\rho(g). %	\rho(g) \overset{!}{=} \rho(eg) = \rho(e)\rho(g)
		\end{align}
		But $eg=g$, so $\rho(g) = \rho(e)\rho(g)$ which is true if and only if $\rho(e) = \id$. Instead taking $h=g^{-1}$:
		\begin{align}
			\rho(g) \rho(g^{-1}) =  \rho(gg^{-1}). %\overset{!}{=} \rho(e) = \id.
		\end{align}
		But $gg^{-1} = e$ and $\rho(e) = \id$, implying $\rho(g^{-1}) = \rho(g)^{-1}$.
	\end{proof}
	
	\begin{note}
		If the vector space $V$ is provided with a basis, then a matrix representation of $G$ may be found. A matrix representation is not canonical, ie it is dependent on the basis chosen for $V$. %Two representations $\rho(g)$ and $\tilde{\rho}(g)$ are similar if there is a transformation $f$ from $\GL(V)$ such that $\tilde{\rho}(g) = f \cdot \rho(g)$ for every $g$ in $G$.
	\end{note}
	
	\begin{note}
		The homomorphism $\rho$, the set of maps $\{\rho(g)\}_{g \in G}$ and the vector space $V$ are interchangeably and abusively called the \textit{representation of $G$}.
	\end{note}
	
	\begin{note}
		Another abuse of notation is that $\rho$ is sometimes dropped and the action of $\rho(g)$ on a vector $\vvec$ is then denoted by $\rho(g)(\vvec) = g\vvec$.
	\end{note}
	
	\subsection{Subrepresentations and irreducibility}
	
		As laid out in \cite{FultonHarris}, we are looking for representations which are said to be ``atomic'' and conversely, which of these indecomposable representations compose any arbitrary representation. %These concepts translate directly from linear algebra. 
		
		\begin{definition}[Analogous concepts in linear algebra and representation theory]\label{def:subreprs}
			These definitions are required to begin the treatment on subrepresentations.		
			\begin{enumerate}
				
				\item The \textbf{degree} of a representation is said to be the dimension of its vector space.
				
				\item A vector space $V$ is said to be \textbf{$G$-invariant}, or fixed under the action of $G$, if for any $g$ in $G$, 
				\begin{align}
					g\vvec \in V,
				\end{align}
				for every $\vvec$ in $V$, or equivalently, $gV \subseteq V$	for every $g$ in $G$.
				
				\item A \textbf{subrepresentation} $W$ of a representation $V$ of $G$ is a $G$-invariant vector subspace $W$ of $V$.
				
				\item A representation is said to be \textbf{irreducible} if it has no subrepresentations, ie it has no proper non-zero invariant subspace.
			\end{enumerate}
		\end{definition}

		\begin{corollary}[Degree 1 representations are irreducible]
			If a representation is of degree 1, its vector space has no subspace other than itself and the zero space, thus it is clearly irreducible.
		\end{corollary}	
		
		Analogously to the existence of orthogonal subspaces in linear algebra (Lemma \ref{lemma:compsubspace}), there are also complementary subrepresentations.
		
		\begin{proposition}\cite[Prop.1.5]{FultonHarris}\label{prop:compsubrepr}
			If $W$ is a subrepresentation of a representation $V$ of $G$, then there is a complementary subspace $W^\perp$ of $V$ such that 
			\begin{enumerate}
				
				\item[i)] $V = W \oplus W^\perp$.
				
				\item[ii)] If $\dim V = n$ and $\dim W = k$, then $\dim W^\perp = n-k$.
				
				\item[iii)] If $(e_i)_{i=1}^n$ is the basis of $V$ and $(e_i)_{i=1}^k$ is the basis of $W$, then $(e_i)_{i=k+1}^n$ is the basis of $W^\perp$.
				
				\item[iv)] $W^\perp$ is also $G$-invariant.
				
			\end{enumerate}
		\end{proposition}
		\begin{proof}
			Items i-iii follow from lemma \ref{lemma:compsubspace}. Let $\vvec \in W^\perp$ and $g \in G$. Then either
			\[
			\begin{cases}
				g\vvec \in W^\perp \text{, and we are done, or }\\
				g\vvec \in W.
			\end{cases}
			\]
			However, the second alternative implies that $\vvec \in W$ since $W$ is $G$-invariant, which is contradictory to the initial assumption that $\vvec \in W^\perp$, therefore $g\vvec \in W$ and $W^\perp$ is $G$-invariant.
		\end{proof}
		
		The results of proposition \ref{prop:compsubrepr} invites the next theorem.
		
		\begin{theorem}[Maschke's theorem]\label{thm:maschke}
			Let $G$ be a finite group and let $V$ be any representation of $G$, then $V$ is composed as a direct sum of a number of not necessarily distinct irreducible subrepresentations $W_i$ of $V$, or in other words,
			\begin{align}
				V = \bigoplus_i W_i.
			\end{align}
			\begin{note}
				The direct sum can be expressed in terms of \underline{distinct} subrepresentations $W_i$ of $V$ and their multiplicities $a_i$ as
				\begin{align}
					V = \bigoplus_i a_i W_i,
				\end{align}
				where $a_i W_i$ is the direct sum of $a_i$ copies of $W_i$.
			\end{note}
		\end{theorem}
		\begin{proof}
			The theorem will be proved using induction on complementary subspaces, using proposition \ref{prop:compsubrepr}. If $V$ itself is irreducible, we are done. If $V$ is not irreducible, then by definition \ref{def:subreprs} there is a subrepresentation $W$ of $V$. Then according to proposition \ref{prop:compsubrepr}, there is a complementary subrepresentation $W^\perp$ to $W$ in $V$ such that $V = W \oplus W^\perp$. If both of $W$ and $W^\perp$ are irreducible, then we are done. If either are not, then use the same method as used on $V$. After recursion on $W$ and $W^\perp$ and since $V$ is a finite vector space, we are done.
		\end{proof}
	
	\subsection{Tensor operations on representations}
	
	In the previous section, we reduced representations into their constituent irreducibles. In this section we will study some methods to create new representations by conversely putting them together using vector space and tensor operations.
	
	Let $V$ and $W$ be representations of $G$ with respective degrees $k$ and $l$ and bases $(\vhat_i)_{i=1}^k$ and $(\what_i)_{i=1}^l$. Then all the following vector spaces are also representations of $G$.
	
	\paragraph{Direct sum of representations.} The direct sum $V \oplus W$ is a representation. By recursion, for a positive integer $n$, the direct sum of $n$ copies of $V$, denoted
	\[
		nV = \underset{n \text{ times}}{\underbrace{V \oplus \cdots \otimes V}}
	\]
	is a representation.
	
	\paragraph{Tensor product of representations.} The tensor product $V \otimes W$ is a representation.
	
	\paragraph{Tensor power of a representation.} By recursion, for a positive integer $n$, the $n$th tensor power of $V$,
	\[
		V^{\otimes n} = \underset{n \text{ times}}{\underbrace{V \otimes \cdots \otimes V}}
	\]
	is a representation. The zeroth tensor power is by definition the ground field, and the first power is the representation itself. 
	
	\paragraph{The symmetric and exterior powers.} The $n$th tensor power $V^{\otimes n}$ has two  subrepresentations: the symmetric powers $\Sym^n V$ and the exterior powers $ \bigwedge^n V$. They are complementary such that $V^{\otimes n} = \Sym^n V \oplus \bigwedge^n V$.
	
	\paragraph{Dual representation.} The dual space of $V$, denoted $V^*$, is a representation. It is the set of all complex transposed vectors\footnote{fysikervarning, $V$ är ett kolumnvektorrum (ket-rum) och $V^*$ är det relaterade radvektorrummet (bra-rum).} of $V$. It is identified with the set of all linear functions\footnote{någon sorts skalär/inre produkt} from $V$ to $\C$.
	
	\paragraph{Set of homomorphisms.} The set of all linear functions between $V$ and $W$, denoted $\Hom(V,W)$ is itself a representation. It is identified with the tensor product $\Hom(V,W) = V^* \otimes W$. In particular, the dual representation is identified with $V^* = \Hom(V,\C)$.
	
	
	
	
	\subsection{Morphismsm of representations}

Let $G$ be a finite group and let $\rho_V$ and $\rho_W$ be two representations of $G$ in the vector spaces $V$ and respectively $W$ such that $\rho_V: G \rightarrow \GL(V)$ and $\rho_W: G \rightarrow \GL(W)$, then a map $\varphi$ from the two representations is a map $\varphi:V \rightarrow W$.

If the map $\varphi$ is such that the diagram
\begin{figure}[hbt!]
	\centering
	\[\begin{tikzcd}
		V \arrow{r}{\varphi} \arrow[swap]{d}{\rho_V(g)} & W \arrow{d}{\rho_W(g)} \\
		V \arrow{r}{\varphi} & W
	\end{tikzcd}\]
	\caption{Morphism of representations}
	\label{fig:commapmorp}
\end{figure}

is commutative for every $g \in G$, it is called a $G$-linear map, or a \emph{morphism of representations}\cite{FultonHarris}. Since this holds for any $\vvec \in V$, we have that
\begin{align}\label{eq:Glinmap}
	\rho_g \circ \varphi = \varphi \circ \rho_g
\end{align}
for any $g \in G$.

The kernel of the map, $\ker \varphi$, defined as $\{\vvec \in V \ | \ \varphi(\vvec) = \0\}$, is a subrepresentation of $V$ since it is a subspace of $V$ \cite{Holst} and $G$-invariant, since clearly for a $\vvec \in \ker \varphi$, $\rho_g \circ \varphi(\vvec) = 0$, then also $\rho_g \circ \varphi(\vvec) = \0$, implying $\rho_g \circ \varphi(\vvec) \in \ker \varphi$.

The image of the map, $\Im \phi$, defined as $$\{ \phi(\vvec) \ | \ \vvec \in V \} = \{ \wvec \in W \ | \ \exists \vvec \in V \text{ s.t. } \varphi(\vvec) = \wvec \},$$ is similarly a subspace of $W$. Let $\wvec \in \Im \varphi$. Thne there exists a $\vvec \in V$ such that $\varphi(\vvec)=\wvec$, then $\rho_g \circ \varphi (\vvec) = \rho_g(\wvec)$, but by Equation \ref{eq:Glinmap}, then also $\varphi \circ \rho_g(\vvec) = \rho_g (\wvec)$, implying $\rho_g(\vvec) \in \Im \varphi$. The image is also $G$-invariant, and thus is a subrepresentation of $W$.

\subsubsection{Subrepresentations as a morphism of representations}

Let $V$ be a representation of a group $G$, and let $W$ be a $G$-invariant subspace, then the projection $\rho: V \rightarrow W$ 



	\subsection{Trivial representation}

	Any group can be represented by a trivial representation of degree 1 defined by 
	\begin{align}
		\rho(g) = 1, \text{ for every } g \in G.
	\end{align}
	It is clearly a homomorphism\footnote{Here I did not state the vector space, should I? ``A vector space over the trivial group/field	''.}. 
	
	For any $G$, and choosing $V = \R$, a map
	\[
	\rho: G \rightarrow \GL(\R)
	\]
	defined by $ \rho(g) = 1$ for every $g$ in $G$, in other words mapping every element in $G$ to the trivial element of $V = \R$. It is trivially a homomorphism since $\rho(e)=1$ and both $\rho(gh)=1$ and $\rho(g)\rho(h)=1$ and is called the \textbf{trivial representation}. It is a representation of degree 1.
	
	%Letting $\mathbf{X}(g) = 1$ for every $g$ in $G$, we have a trivial representation of degree 1. This is trivially a homomorphism and thus a representation.
	
	Any group has a trivial degree 1 representation, but there are many trivial representations of larger degree, for example choosing $V = \R^3$, we define the representation by $\rho(g) = \1_3$, the three-by-three identity matrix. However, this is not an irreducible representation as $\R^3 = \R \oplus \R \oplus \R = 3\R$ and
	\[
	\begin{pmatrix} 1&0&0\\0&1&0\\0&0&1 \end{pmatrix} = (1)\oplus(1)\oplus(1).
	\]
	
	For the symmetric group $S_n$, the \textbf{trivial representation} is the one-dimensional representation that maps every permutation to 1.
	\[
	\rho(\sigma) = \1, \ \forall \sigma \in S_n
	\]
	It is a representation because in particular it maps $e$ to 1 and for any $\sigma$, $\tau$ in $S_n$,
	\[
	\rho(\sigma \circ \tau) = 1 = 1 \cdot 1 = \rho(\sigma) \cdot \rho(\tau)
	\]
	it is a homomorphism.

\subsection{Alternating representation of $S_n$}

	Choosing $G$ to be the symmetric group of degree $n$, denoted $S_n$, we can find another degree 1 representation by studying the sign, or parity, of the elements of $S_n$. For any two permutations $\sigma$ and $\tau$ of $S_n$ with respective sign $(-1)^s$ and $(-1)^t$, the sign of their composition is 
	\begin{align}
		\sgn(\sigma \circ \tau) = (-1)^s \cdot (-1)^t = (-1)^{st} = \sgn(\sigma) \cdot \sgn(\tau),
	\end{align}
	so clearly this sign function is a homomorphism from $S_n$ to a one-dimensional vector space over the two elements\footnote{The vector space is over the field of two elements, $\F_2$.} $+1$ and $-1$, hence the $\sgn$ function is representation of degree 1 and is called the alternating representation, or alternatively the sign or signature representation.
	
	The symmetric group has another degree 1 representation. Every element $\sigma$ of the symmetric group of degree $n$, denoted $S_n$, has a sign, or parity, denoted $\sgn(\sigma)$ defined by the number of transpositions\footnote{2-cycles} required to construct it. Letting $k$ be that number, then the parity or sign of $\sigma$ is $\sgn(\sigma)=(-1)^k$, so an \textbf{even} permutation has sign $+1$ and an \textbf{odd} permutation has sign $-1$. Now, with $G = S_n$ and choosing $V=\F_2$, this $\sgn$ function
	\[
	\sgn: S_n \rightarrow \{-1, 1\}
	\]
	defined as above is also a homomorphism since $\rho(e) = 1$ and for two arbitrary elements $\sigma$ and $\tau$ of $S_n$ with respective sign $(-1)^s$ and $(-1)^t$ we have 
	\[
	\rho(\sigma)\rho(\tau) = (-1)^s(-1)^t = (-1)^{st} = \rho(\sigma\tau).
	\] 
	This $\sgn$ function is called the \textbf{alternating representation}\footnote{In some literature called the sign or signature representation.} of the symmetric group.
	
	Recall that all of $S_n$ can be generated by its subset of adjacent transpositions $\{(12), (23), \dots, (n,1)\}$, so any permutation in $S_n$ can be written as a product of transpositions. If the number of transpositions required to compose the permutation $\sigma$ is $k$, then the sign of $\sigma$ is defined as
	\[
	\sgn(\sigma) = (-1)^k.
	\]
	So for an even(odd) $k$, the permutation is called even(odd). The $\sgn$ function arises naturally from the symmetric group and is called the \textbf{alternating\footnote{It is sometimes called the sign or signature representation} representation}. It is a representation because it maps $e$ to 1 and a composition $\sigma \circ \tau$ has the sign $\sgn(\sigma\tau) =\sgn(\sigma)\sgn(\tau)$, hence it is a homomorphism.
	
	

\subsection{Permutation representation}

	We could let the group $G$ act on an appropriately chosen set $X$, and study its action $f$ by permutation of the elements of $X$. This action is defined by any $g$ in $G$ and $x$ in $X$ as
	\begin{table}[hbt!]\centering\begin{tabular}{c c c c}
			$f:$ & $G \times X$    & $\rightarrow$ & $X$ \\
			&\rotatebox[origin=c]{90}{$\in$}&&\rotatebox[origin=c]{90}{$\in$} \\
			$f:$ & $(g,x)$ & $\mapsto$     & $gx$.
	\end{tabular}\end{table}
	
	Further more, define a vector space $V$ spanned by the orthogonal basis $\mathcal{B} = \{\ehat_x\}_x$ for every $x$ in $X$. Then we have a homomorphism
	\begin{align}
		\rho: G \rightarrow \GL(V)
	\end{align}
	defined by its action on the basis $\mathcal{B}$ by, for any $g$ in $G$,
	\begin{align}
		\rho(g): \ehat_x \mapsto g\ehat_x = \ehat_{gx}.
	\end{align}
	or on an arbitrary vector $v$ in $V$ as
	\begin{align}
		\rho(g): \vvec \mapsto g\vvec.
	\end{align}
	
	\textit{Show that it is a homomorphism!}
	
	\subsubsection{$S_n$ acting on $\N_n$}
	The action is defined as $S_n \times \N_n$ by $(\sigma,i) \mapsto \sigma(i)$. Choosing a basis $\mathcal{B} = \{\ehat_i\}_{i=1}^n$ we span a vector space $V$ of dimension $\dim V = n$. We then define a homomorphism $\rho$ from $S_n$ to $\GL(V)$ by
	\begin{align}
		\rho(\sigma): \ehat_i \mapsto \ehat_{\sigma(i)},
	\end{align}
	that is by permutating a basis vector $\ehat_i$ to another one $\ehat_j = \ehat_{\sigma(i)}$.
	
	\textit{action of $\rho(\sigma)$ on any vector??}
	
	The corresponding matrix representations are the permutation matrices
	\begin{align}
		X(\sigma) = (x_{ij})_{n \times n}, \text{ where } x_{ij} = \delta_{j,\sigma(i)} = \begin{cases}
			1 \text{ if } j = \sigma(i),\\
			0 \text{ otherwise.}
		\end{cases}
	\end{align}
	
	\begin{example}[$S_n$ for $n = 2,3$]
		
		The symmetric group of degree 2 has two elements, $S_2 = \{(1), (12)\}$ and their matrix representations in $\mathbf{\ehat_1},\mathbf{\ehat_2}$-space are
		\begin{table}[hbt!]
			\centering
			\caption{Matrix representations of $S_2$}
			\begin{tabular}{c | c c}
				$\sigma$    & $(1)$ & $(12)$ \\
				\hline
				$X(\sigma)$ & $\begin{pmatrix}
					1 & 0 \\ 0 & 1
				\end{pmatrix}$ & $\begin{pmatrix}
					0 & 1 \\ 1 & 0
				\end{pmatrix}$
			\end{tabular}
		\end{table}
		
		Likewise, representations of $S_3$ follows:
		
		\begin{table}[hbt!]
			\centering
			\caption{Matrix representations of $S_3$}
			\begin{tabular}{c | c c c c c c}
				$\sigma$  &$(1)$& $(12)$&$(13)$&$(23)$&$(123)$&$(132)$ \\
				\hline
				$X(\sigma)$ & 
				$\left(\begin{smallmatrix}
					1 & 0 & 0 \\
					0 & 1 & 0 \\
					0 & 0 & 1
				\end{smallmatrix}\right)$ &
				$\left(\begin{smallmatrix}
					0 & 1 & 0 \\
					1 & 0 & 0 \\
					0 & 0 & 1
				\end{smallmatrix}\right)$ &
				$\left(\begin{smallmatrix}
					0 & 0 & 1 \\
					0 & 1 & 0 \\
					1 & 0 & 0
				\end{smallmatrix}\right)$ &
				$\left(\begin{smallmatrix}
					1 & 0 & 0 \\
					0 & 0 & 1 \\
					0 & 1 & 0
				\end{smallmatrix}\right)$ &
				$\left(\begin{smallmatrix}
					0 & 0 & 1 \\
					1 & 0 & 0 \\
					0 & 1 & 0
				\end{smallmatrix}\right)$ &
				$\left(\begin{smallmatrix}
					0 & 1 & 0 \\
					0 & 0 & 1 \\
					1 & 0 & 0
				\end{smallmatrix}\right)$
			\end{tabular}
		\end{table}
		
		%			$$S_3 = \{(1), (12), (13), (23), (123), (132)\}$$
		%			$$S_4 = \{(1), 
		%						(12), (13), (14), (23), (24), (34), 
		%						(12)(34), (13)(24), (14)(23), 
		%						(123), (124), (134), (132), (142), (143), (234), (243),
		%						(1234), (1243), (1324), (1342), (1423), (1432)
		%						\}$$
	\end{example}
	
	Let the group $G$ act on an appropriately chosen set $A$ by
	\[
	\phi : G \times A \rightarrow A
	\] % could be replaced by by ``permutating its elements''
	defined by $\phi: g \mapsto ga$. After constructing a vector space\marginnote{What is the dim?} $V$ with a basis $\mathcal{B} = ( \ebas_a )_{a \in A}$, we can define a representation
	\[
	\rho: G \rightarrow \GL(V)
	\]
	defined by
	\[
	\rho(g)(\ebas_a) = \ebas_{ga}.
	\]
	This is a homomorphism since for any $a$ in $A$, $\rho(e)$ acts as a identity map by 
	\[
	\rho(e)(\ebas_a) = \ebas_{ea} = \ebas_a
	\]
	and also for any $g$ and $h$ in $G$, 
	\[
	\rho(g)\rho(h)(\ebas_a) = \rho(g)(\ebas_{ha}) = \ebas_{gha} = \rho(gh)(\ebas_a).
	\]
	This homomorphism is called the \textbf{permutation representation}.
	
	Instead of letting $S_n$ act naturally on the set $\{1, 2, \dots, n\}$, let $S_n$ act on the coordinates of points in $\C^n$ by permuting the basis vectors $(\hat{e}_i)_{i=1}^n$ by $\sigma(\hat{e}_i) = \hat{e}_{\sigma(i)}$, or in other words, $\sigma$ permutes the $i$th to the $\sigma(i)$th axis. Then we have a representation
	\[
	\rho: S_n \longrightarrow \GL_n
	\]
	that is defined by 
	\[
	\rho(\sigma) = \left(\delta_{i, \sigma(i)}\right)_{n \times n},
	\]
	where $\delta_{ij}$ is the Kronecker delta. In other words, every $\rho(\sigma)$ is the unit matrix where the $i$th row is moved to the $\sigma(i)$th row. It is a representation since it maps $e$ to $\1_{n  \times n}$ and it is a homomorphism by matrix multiplication.
	
	\begin{example}[Action of $S_3$ on coordinates in $\R^3$.]
		Letting $S_3$ act on coordinates in space by permuting the axes we obtain:
		\begin{table}[hbt!]
			\centering
			\begin{tabular}{r r r}
				$\rho(e) = 
				\begin{pmatrix}
					1 & 0 & 0 \\
					0 & 1 & 0 \\
					0 & 0 & 1
				\end{pmatrix}$ & 
				$\rho(123) = 
				\begin{pmatrix}
					0 & 0 & 1 \\
					1 & 0 & 0 \\
					0 & 1 & 0
				\end{pmatrix}$ & 
				$\rho(132) = 
				\begin{pmatrix}
					0 & 1 & 0 \\
					0 & 0 & 1 \\
					1 & 0 & 0
				\end{pmatrix}$ \\ & & \\
				$\rho(12) = 
				\begin{pmatrix}
					0 & 1 & 0 \\
					1 & 0 & 0 \\
					0 & 0 & 1
				\end{pmatrix}$ &
				$\rho(13) = 
				\begin{pmatrix}
					0 & 0 & 1 \\
					0 & 1 & 0 \\
					1 & 0 & 0
				\end{pmatrix}$ &
				$\rho(23) = 
				\begin{pmatrix}
					1 & 0 & 0 \\
					0 & 0 & 1 \\
					0 & 1 & 0
				\end{pmatrix}$
			\end{tabular}
		\end{table}
	\end{example}
	
	The question is, is this an irreducible representation? It is not since a representation is irreducible if it has no non-zero proper invariant subspace. Coordinates of points in $\R^3$ along the line from the origin through $(1,1,1)$ is invariant, consequently the subspace $$\left\{ (x,y,z) \in \R^3 \ \middle\vert \ x+y+z=0 \right\}$$ is $S_3$-invariant and is irreducible.

\subsection{Standard representation}

	\textit{The subspace found when studying $S_3$ - generalize to $S_n$.}

\subsection{Regular representation}

	Instead of letting the group $G$ act on some arbitrary set by permutation, we could let it act on itself. This action is defined as
	\[
	G \times G \rightarrow G
	\]
	which for every $g$ in $G$ is given by 
	\[
	(g,h) \mapsto gh,
	\]
	for a $h$ in $G$. The corresponding vector space $V$ is spanned by basis vectors $\{\ghat\}_{g \in G}$ constructed from the members of $G$. A representation of $G$ in $V$ is then
	\[
	\rho: G \rightarrow \GL(V)
	\]
	defined by
	\[
	\rho(g)(\hhat) = g\hhat.
	\]
	
	Instead of any general set, we can let $G$ act on itself. Similarly as above, we construct a vector space with basis $( \ebas_g )_{g \in G}$ and define a representation
	\[
	\rho: G \rightarrow \GL(V)
	\]
	defined by
	\[
	\rho(g)(\ebas_h) = \ebas_{gh}.
	\]
	To show that it is a homomorphism, one could follow the same steps as above. This is called the \textbf{regular representation}.
	
	\subsubsection{The regular representation of cyclic group $Z_n$}
	
		Recall that the cyclic group of order $n$ is generated by an element $g$ of order\footnote{The order of an element $g$ is the smallest possible integer $n$ such that $g^n = e$.} $n$ if $Z_n = \{e, g, g^2, \dots, g^{n-1} \}$. Sometimes we denote the cyclic group as $Z_n = \langle g \rangle$ if $g$ is such a generator. A cyclic group may have multiple generators that generate the group separately, for example the integers modulo 12, $\Z/12\Z \cong Z_{12}$ has 1, 5, 7 and 11 as generators. 
		
		Consider a matrix representation of $Z_n$ as a homomorphism
		\[
		\rho: Z_n \longrightarrow \C
		\]
		defined by $\rho(e) = 1$ and $\rho(g^a g^b) = \rho(g^a)\rho(g^b)$. Since $g^n = e$ and $\rho(g^n) = 1$, then $\rho(g)$ must be mapped to an $n$th root of unity. Denote by $\omega = e^{2\pi i/n}$ the first non-real $n$th root of unity. So for $Z_n$, there are $n$ one-dimensional representations $\rho_m: Z_n \rightarrow \C$ that maps $g^k$ to $\omega^{km}$.
		
		\begin{example}[$Z_3$]
			The three third roots of unity are $1$, $\omega = \frac{-1+i\sqrt{3}}{2}$ and $w^2 = \frac{-1-i\sqrt{3}}{2}$, so the three representations of $Z_3$ are presented in table \ref{table:Z3}.
			
			\begin{table}[hbt!]
				\centering
				\caption{Three representations of $Z_3$.}
				\begin{tabular}{c | c c c}
					\label{table:Z3}
					$Z_3$ & $e$ & $g$        & $g^2$      \\ \hline
					$\rho_0$          & 1   & 1          & 1          \\
					$\rho_1$          & 1   & $\omega$   & $\omega^2$ \\
					$\rho_2$          & 1   & $\omega^2$ & $\omega$
				\end{tabular}
			\end{table}
		\end{example}
		
		\begin{example}[$Z_5$]
			The five fifth roots of unity are $e^{\frac{2\pi im}{5}}$, so five representations of $Z_5$ are presented in table \ref{table:Z5}.
			
			\begin{table}[hbt!]
				\centering
				\caption{Five representations of $Z_5$. $\omega = e^{2 \pi i/5}$.}
				\begin{tabular}{c | c c c c c}
					\label{table:Z5}
					$Z_5$ & $e$ & $g$        & $g^2$      & $g^3$      & $g^4$      \\ \hline
					$\rho_0$            & 1   & 1          & 1          & 1          & 1          \\
					$\rho_1$            & 1   & $\omega^1$ & $\omega^2$ & $\omega^3$ & $\omega^4$ \\
					$\rho_2$            & 1   & $\omega^2$ & $\omega^4$ & $\omega^1$ & $\omega^3$ \\
					$\rho_3$            & 1   & $\omega^3$ & $\omega^1$ & $\omega^4$ & $\omega^2$ \\
					$\rho_4$            & 1   & $\omega^4$ & $\omega^3$ & $\omega^2$ & $\omega^1$
				\end{tabular}
			\end{table}
		\end{example}
		
		
	\subsection{Dihedral group $\D_{2n}$}
	
		The dihedral group of order 8, degree 4, $\D_8$ is generated by rotation $r$ and mirroring $s$ such that
		
		\begin{align}
			\D_8 = \left\langle r, s \ \middle\vert \ r^4 = s^2 = e, rs = sr^-1 \right\rangle
		\end{align}
		
		It is a 2-Sylow subgroup of $S_4$.
		
		It has a degree 2 representation generated by
		
		\begin{align}
			R = \begin{pmatrix}
				0 & -1 \\
				1 & 0 
			\end{pmatrix} \ \text{ and } \ S = \begin{pmatrix}
			1 & 0 \\
			0 & -1 
			\end{pmatrix}.
		\end{align}
		
		One could check that $R^4 = S^2 = \1_{2 \times 2}$ and $RS = SR^-1$.
		
		$\D_8$ has five conjugacy classes, $[e], [r^2], [r, r^3], [s, r^2s], [rs, r^3s]$.
		
		\begin{table}[hbt!]\caption{Character table for $\D_8$}\centering
			\begin{tabular}{c | c c c c c }
				 $[g]$  & $[e]$ & $[r^2]$ & $[r, r^3]$ & $[s, r^2s]$ & $[rs, r^3s]$ \\
				$|[g]|$ & 1     & 1       & 2          & 2           & 2            \\ \hline
				 Triv   & 1     & 1       & 1          & 1           & 1            \\
				 2-dim  & 2     & -2      & 0          & 0           & 0            \\
				     ?  &       &         &            &             &              \\
				     ?  &       &         &            &             &              \\
				     ?  &       &         &            &             &
			\end{tabular}
		\end{table}
		
		
			
			
%	\clearpage{\thispagestyle{empty}}
\section{Representation Theory 1}

This section is based on the treatment on representation theory from \cite{FultonHarris}.

For the rest of this thesis, we are considering finite groups and finite-dimensional vector space over $\C$. Let $G$ be such a group and $V$ such a vector space. The group of all automorphism\footnote{This is equivalent to the set of invertible $n \times n$ matrices with entries from $\C$.} of $V$ is denoted by $\GL(V)$.

\begin{definition}[Representation]
	A representation of $G$ on the vector space $V$ is a homomorphism
	\[
	\rho: G \rightarrow \GL(V) 
	\]
	defined by taking each element $g$ in $G$ and mapping it to an element of $\GL(V)$. Equivalently it is a mapping $\rho(g): V \rightarrow V$. 
\end{definition}
\begin{remark}
	Since $\rho$ is a homomorphism, 
	\begin{enumerate}[i)]
		\item $\rho(e) = \1$, or $\rho(e)(v)=v$ for every $v$ in $V$.
		\item For any $g$ and $g'$ in $G$, $\rho(gg') = \rho(g)\rho(g')$.
		\item For any $g$ and its inverse $g^{-1}$ in $G$, $\rho(g^{-1}) = \rho(g)^{-1}$.
	\end{enumerate}
\end{remark}

In other words, the matrix $\rho(g)$ acts acts as a linear map on a vector $v$ in $V$ and maps it to the vector $\rho(g)(v)$ in $V$. Sometimes the vector space $V$ is itself shorthandedly called the representation and $\rho(g)(v)$ is abbreviated to $gv$.

%	\begin{example}[Trivial representation]
	%		By choosing to represent every $g \in G$ by the identity matrix, we get the trivial representation of $G$. Clearly this is a homomorphism since for $g,h \in G$, $\rho(gh) = 1 = \rho(g) \rho(h)$.
	%	\end{example}

\begin{definition}[Subrepresentation]
	A subrepresentation $W$ of a representation $V$ of $G$ is a subspace $W$ of $V$ invariant under $G$.
\end{definition}

%	\begin{example}[Permutation representation]\label{example:C3}
	%		Studying the action of $S_3$ on coordinates in the complex space $\C^3$, it is natural to consider $\C^3$ as a permutation representation. However, $\C^3$ is not invariant under $S_3$, eg. the line through $(1,1,1)$ and origin is fixed under $S_3$, so the subspace $W = \{ (z_1, z_2, z_3) \in \C^3 \mid z_1+z_2+z_3 = 0 \}$ is invariant under $S_3$.
	%	\end{example}

\begin{definition}[Irreducible representation]
	A representation $V$ is irreducible if it has no proper non-zero invariant subspace.
\end{definition}

%	\begin{example}
	%		The permutation representation $V=\C^3$ is not irredicible since $W$ from example \eqref{example:C3} is a proper non-zero invariant subspace of $V$. However, this subrepresentation $W$ is irreducible. 
	%	\end{example}

\subsection{Tensor operations on representations}

Let $V$ and $W$ be two representations, then their direct sum $V \oplus W$, their tensor product $V \otimes W$ are representations.

The $n$th tensor power of $V$, \[V^{\otimes n} = \underset{n \text{ times}}{\underbrace{V \otimes \cdots \otimes V}}\] is also a representation with the symmetric powers $\Sym^n V$ and the exterior powers $ \bigwedge^n V$ as subrepresentations. Further more, the dual space $V^* = \Hom(V,\C)$ of $V$ and more generally $\Hom(V,W) = V^* \otimes W$ are also representations. Recall that $\Hom(V,W)$ is the set of all linear functions (functionals??) from $V$ to $W$ and in particular $\Hom(V, \C)$ ``is some sort of scalar product''. Cf. quantum physics, if $v$ is a column vector from a ket-space $V$, then $u^\dagger$ is a row vector from the corresponding bra-space $V^*$.

%\clearpage{\thispagestyle{empty}}
%\section{Character Theory}

%\clearpage{\thispagestyle{empty}}
%\section{Representations of some Finite Groups}

\subsection{The symmetric group $S_n$}

(Purpose: Recall the symmetric group and ultimately show that the character function is a class function (=invariant in conjugacy class).)

This section is based on the first chapter of \cite{Sagan}.

Recall that the symmetric group $S_n$ is the set of bijections from $\{ 1, 2, \dots, n\}$ to itself. Its elements are called \textbf{permutations} and it is a group under composition of permutations. A permutation $\sigma$ in $S_n$ can be represented 
%	in two-line notation
%	\[ 
%	\sigma = 
%	\begin{pmatrix}
	%		1 & 2 & \cdots & n \\
	%		\sigma(1) & \sigma(2) & \cdots & \sigma(n)
	%	\end{pmatrix},
%	\]
%	or more commonly 
in cycle notation as a composition of disjoint cycles
\[
\sigma = \big( i, \sigma(i), \sigma^2(i), \cdots, \sigma^{-1}(i) \big) \circ \big( j, \sigma(j), \sigma^2(j), \cdots, \sigma^{-1}(j) \big) \circ \cdots.
\]
which carries information on how $\{1,2,\dots,n\}$ is permutated. \begin{example}
	For example, the element $\tau$ of $S_5$ that is defined by 
	\[
	\sigma(1) = 2 \quad \sigma(2) = 1 \quad \sigma(3) = 3 \quad \sigma(4) = 5 \quad \sigma(5) = 4
	\]
	can be represented as 
	\[ 
	\tau = 
	\begin{pmatrix}
		1 & 2 & 3 & 4 & 5 \\
		2 & 1 & 3 & 5 & 4
	\end{pmatrix}
	\]
	or as $\tau = (12)(3)(45) = (12)(45)$, consisting of a 1-cycle that is usually omitted and two 2-cycles. A 2-cycle $(ij)$ is called a \textbf{transposition}.
\end{example}

The \textbf{cycle type} of $\sigma$ in $S_n$ is an $n$-tuple $(k^{m_k})_{k=1}^n$ where $m_k$ is the number of $k$-cycles in $\sigma$. The cycle type of the example above is $\text{Type}(\tau) = (1^1,2^2,3^0,4^0,5^0) = (1^1,2^2)$, where the cycles of multiplicity 0 are omitted.

Another way to describe a permutation in $S_n$ is to compare it to an integer partition of $n$. The integer partition of $n$ associated with $\sigma$ in $S_n$ is denoted by $\lambda(\sigma) = ( \lambda_1, \lambda_2, \dots, \lambda_l)$ with $l \leq n$, $\lambda_i \leq \lambda_{i+1}$ and $\sum_{i=1}^l \lambda_l = n$. The permutation $\tau$ from earlier corresponds to the partition $\lambda(\tau) = (2, 2, 1)$. 
\ytableausetup{smalltableaux}
A third equivalent way is to draw the Young diagram of the associated partition, for example the Young diagram of $\tau$ from above is 
\begin{table}[hbt!]
	\centering
	\ydiagram{2,2,1}.
\end{table}

\begin{example}[$S_3$]
	Table \ref{table:S3} present the elements of $S_3$.	
	\begin{table}[hbt!]
		\centering
		\caption{The elements of $S_3$.}
		\begin{tabular}{r | c c c c c c}
			\label{table:S3}
			$S_3$ & e                & (12)           & (13)           & (23)           & (123)        & (132)        \\ \hline
			Type & $(1^3)$          & $(1^1,2^1)$    & $(1^1,2^1)$    & $(1^1,2^1)$    & $(3^1)$      & $(3^1)$      \\
			$\lambda$ & $(1,1,1)$        & $(2,1)$        & $(2,1)$        & $(2,1)$        & $(3)$        & $(3)$        \\
			Young & \ydiagram{1,1,1} & \ydiagram{2,1} & \ydiagram{2,1} & \ydiagram{2,1} & \ydiagram{3} & \ydiagram{3}
		\end{tabular}
	\end{table}
\end{example}

\subsubsection{Conjugacy classes of $S_n$}

For larger $n$, the number of elements increase very quickly since the size of $S_n$ is $n!$, however as we'll see, we can instead study the conjugacy classes of $S_n$. Recall that two elements $g$ and $g'$ in $G$ are said to be \textbf{conjugate} if there exists an element $h$ in $G$ such that $hgh = g'$. ``Being conjugate'' in a group is an equivalence relation, so the equivalence classes (called \textbf{conjugacy classes}) partition $G$ into disjoint subsets. Denote the conjugacy class of an element $g$ in $G$ as $[g]$. If another element $g'$ is conjugate to $g$, then they share conjugacy class $[g] = [g']$ and both $g$ and $g'$ are said to be \textbf{representatives} of their common conjugacy class.

The size of the conjugacy class $[g]$ can be calculated with the centralizer of $g$ in $G$, defined by 
\[
\Cent(g) = \left\{ h \in G \ \middle\vert \ hgh^{-1} = g \right\}
\]
and by the orbit-stabilizer theorem\cite[Thm X.X]{DummitFoote}, the relationship between $\Cent(g)$ and the elements of $[g]$ is
\[
|[g]| = \frac{|G|}{|\Cent(g)|}
\]
Returning to the symmetric group, two permutations $\sigma$ and $\tau$ share conjugacy class if and only if they are of the same cycle type\cite{DF}. Since the cycle type of a permutation was linked to a integer partition of the degree of the symmetric group, there are as many conjugacy classes of $S_n$ as there are integer partitions of $n$. Eg. there are 3 conjugacy classes in $S_3$ and 7 in $S_5$. Also, if the cycle type of a permutation $\sigma$ is $(k^{m_k})_{k=1}^n$, then the size of its centralizer is $|\Cent(\sigma)| = \prod_{k=1}^{n} k^{m_k}{m_k}!$\cite{Sagan}. 


\begin{example}[$S_5$]
	The conjugacy classes of $S_5$, along with their sizes and their types are presented in table \ref{table:S5}.
	\begin{table}[hbt!]
		\centering
		\caption{The conjugacy classes of $S_5$.}
		\begin{tabular}{r | c c c c c c c}
			\label{table:S5}
			$S_5$ & [e]                  & [12]               & [(12)(34)]       & [123]            & [1234]         & [(12)(345)]    & [12345]      \\ \hline
			$|\Cent(\sigma)|$ & $5!$                 & 12                 & 8                & 6                & 6              & 4              & 5            \\
			$|[\sigma]|$ & 1                    & 10                 & 15               & 20               & 20             & 30             & 24           \\
			Type & $(1^5)$              & $(1^3,2^1)$        & $(1^1,2^2)$      & $(1^2,3^1)$      & $(1^1,4^1)$    & $(2^1,3^1)$    & $(5^1)$      \\
			$\lambda $ & (1,1,1,1,1)          & (2,1,1,1)          & (2,2,1)          & (3,1,1)          & (4,1)          & (3,2)          & (5)          \\
			Young & \ydiagram{1,1,1,1,1} & \ydiagram{2,1,1,1} & \ydiagram{2,2,1} & \ydiagram{3,1,1} & \ydiagram{4,1} & \ydiagram{3,2} & \ydiagram{5}
\end{tabular}\end{table}\end{example}

\subsubsection{Trivial and alternating representation of $S_n$}

For the symmetric group $S_n$, the \textbf{trivial representation} is the one-dimensional representation that maps every permutation to 1.
\[
\rho(\sigma) = \1, \ \forall \sigma \in S_n
\]
It is a representation because in particular it maps $e$ to 1 and for any $\sigma$, $\tau$ in $S_n$,
\[
\rho(\sigma \circ \tau) = 1 = 1 \cdot 1 = \rho(\sigma) \cdot \rho(\tau)
\]
it is a homomorphism.

Recall that all of $S_n$ can be generated by its subset of adjacent transpositions $\{(12), (23), \dots, (n,1)\}$, so any permutation in $S_n$ can be written as a product of transpositions. If the number of transpositions required to compose the permutation $\sigma$ is $k$, then the sign of $\sigma$ is defined as
\[
\sgn(\sigma) = (-1)^k.
\]
So for an even(odd) $k$, the permutation is called even(odd). The $\sgn$ function arises naturally from the symmetric group and is called the \textbf{alternating\footnote{It is sometimes called the sign or signature representation} representation}. It is a representation because it maps $e$ to 1 and a composition $\sigma \circ \tau$ has the sign $\sgn(\sigma\tau) =\sgn(\sigma)\sgn(\tau)$, hence it is a homomorphism.

\subsubsection{Permutation of coordinates}

Instead of letting $S_n$ act naturally on the set $\{1, 2, \dots, n\}$, let $S_n$ act on the coordinates of points in $\C^n$ by permuting the basis vectors $(\hat{e}_i)_{i=1}^n$ by $\sigma(\hat{e}_i) = \hat{e}_{\sigma(i)}$, or in other words, $\sigma$ permutes the $i$th to the $\sigma(i)$th axis. Then we have a representation
\[
\rho: S_n \longrightarrow \GL_n
\]
that is defined by 
\[
\rho(\sigma) = \left(\delta_{i, \sigma(i)}\right)_{n \times n},
\]
where $\delta_{ij}$ is the Kronecker delta. In other words, every $\rho(\sigma)$ is the unit matrix where the $i$th row is moved to the $\sigma(i)$th row. It is a representation since it maps $e$ to $\1_{n  \times n}$ and it is a homomorphism by matrix multiplication.

\begin{example}[Action of $S_3$ on coordinates in $\R^3$.]
	Letting $S_3$ act on coordinates in space by permuting the axes we obtain:
	\begin{table}[hbt!]
		\centering
		\begin{tabular}{r r r}
			$\rho(e) = 
			\begin{pmatrix}
				1 & 0 & 0 \\
				0 & 1 & 0 \\
				0 & 0 & 1
			\end{pmatrix}$ & 
			$\rho(123) = 
			\begin{pmatrix}
				0 & 0 & 1 \\
				1 & 0 & 0 \\
				0 & 1 & 0
			\end{pmatrix}$ & 
			$\rho(132) = 
			\begin{pmatrix}
				0 & 1 & 0 \\
				0 & 0 & 1 \\
				1 & 0 & 0
			\end{pmatrix}$ \\ & & \\
			$\rho(12) = 
			\begin{pmatrix}
				0 & 1 & 0 \\
				1 & 0 & 0 \\
				0 & 0 & 1
			\end{pmatrix}$ &
			$\rho(13) = 
			\begin{pmatrix}
				0 & 0 & 1 \\
				0 & 1 & 0 \\
				1 & 0 & 0
			\end{pmatrix}$ &
			$\rho(23) = 
			\begin{pmatrix}
				1 & 0 & 0 \\
				0 & 0 & 1 \\
				0 & 1 & 0
			\end{pmatrix}$
		\end{tabular}
	\end{table}
\end{example}

The question is, is this an irreducible representation? It is not since a representation is irreducible if it has no non-zero proper invariant subspace. Coordinates of points in $\R^3$ along the line from the origin through $(1,1,1)$ is invariant, consequently the subspace $$\left\{ (x,y,z) \in \R^3 \ \middle\vert \ x+y+z=0 \right\}$$ is $S_3$-invariant and is irreducible.

\subsection{The cyclic group $Z_n$}

Recall that the cyclic group of order $n$ is generated by an element $g$ of order\footnote{The order of an element $g$ is the smallest possible integer $n$ such that $g^n = e$.} $n$ if $Z_n = \{e, g, g^2, \dots, g^{n-1} \}$. Sometimes we denote the cyclic group as $Z_n = \langle g \rangle$ if $g$ is such a generator. A cyclic group may have multiple generators that generate the group separately, for example the integers modulo 12, $\Z/12\Z \cong Z_{12}$ has 1, 5, 7 and 11 as generators. 

Consider a matrix representation of $Z_n$ as a homomorphism
\[
\rho: Z_n \longrightarrow \C
\]
defined by $\rho(e) = 1$ and $\rho(g^a g^b) = \rho(g^a)\rho(g^b)$. Since $g^n = e$ and $\rho(g^n) = 1$, then $\rho(g)$ must be mapped to an $n$th root of unity. Denote by $\omega = e^{2\pi i/n}$ the first non-real $n$th root of unity. So for $Z_n$, there are $n$ one-dimensional representations $\rho_m: Z_n \rightarrow \C$ that maps $g^k$ to $\omega^{km}$.

\begin{example}[$Z_3$]
	The three third roots of unity are $1$, $\omega = \frac{-1+i\sqrt{3}}{2}$ and $w^2 = \frac{-1-i\sqrt{3}}{2}$, so the three representations of $Z_3$ are presented in table \ref{table:Z3}.
	
	\begin{table}[hbt!]
		\centering
		\caption{Three representations of $Z_3$.}
		\begin{tabular}{c | c c c}
			\label{table:Z3}
			$Z_3$ & $e$ & $g$        & $g^2$      \\ \hline
			$\rho_0$          & 1   & 1          & 1          \\
			$\rho_1$          & 1   & $\omega$   & $\omega^2$ \\
			$\rho_2$          & 1   & $\omega^2$ & $\omega$
		\end{tabular}
	\end{table}
\end{example}

\begin{example}[$Z_5$]
	The five fifth roots of unity are $e^{\frac{2\pi im}{5}}$, so five representations of $Z_5$ are presented in table \ref{table:Z5}.
	
	\begin{table}[hbt!]
		\centering
		\caption{Five representations of $Z_5$. $\omega = e^{2 \pi i/5}$.}
		\begin{tabular}{c | c c c c c}
			\label{table:Z5}
			$Z_5$ & $e$ & $g$        & $g^2$      & $g^3$      & $g^4$      \\ \hline
			$\rho_0$            & 1   & 1          & 1          & 1          & 1          \\
			$\rho_1$            & 1   & $\omega^1$ & $\omega^2$ & $\omega^3$ & $\omega^4$ \\
			$\rho_2$            & 1   & $\omega^2$ & $\omega^4$ & $\omega^1$ & $\omega^3$ \\
			$\rho_3$            & 1   & $\omega^3$ & $\omega^1$ & $\omega^4$ & $\omega^2$ \\
			$\rho_4$            & 1   & $\omega^4$ & $\omega^3$ & $\omega^2$ & $\omega^1$
		\end{tabular}
	\end{table}
\end{example}
 % RENSAD
%		\clearpage{\thispagestyle{empty}}
\section{Representation Theory 2}

The layout of this section is based on \cite{Serre}.

Let $V$ be a finite vector space over $\C$ with dimension $\dim V = n$. The group of automorphisms of $V$ is denoted $\GL(V)$ and any element $f$ therefrom is a linear map $f: V \rightarrow V$. If a basis is provided for $V$ then $\GL(V)$ is isomorphic to $\GL_n(\C)$, the set of invertible $n \times n$-matrices, and then $f$ is such a matrix, where again $n = \dim V$\cite[18.1]{DummitFoote}.

Let $G$ be a finite group of degree $|G| = n$ written multiplicatively. 

\begin{definition}[Representation]
	A representation of a finite group $G$ in the vector space $V$ is a homomorphism 
	\[
	\rho: G \rightarrow \GL(V)
	\]
	that assign every element of $G$ to an element of $\GL(V)$. If $V$ has a basis, there is for every $g$ in $G$ a matrix representation
	\[
		\mathbf{X}(g) = (x_{ij})_{n \times n}.
	\]
	\begin{note}
		As a shorthand, $V$ itself is called the representation.
	\end{note}
\end{definition}

As a consequence of the homomorphism of $\rho$, $\rho(e) = \1_n$ (the identity matrix) for the identity element $e$ of $G$ and $\rho(gh) = \rho(g)\rho(h)$ and $\rho(g)^{-1} = \rho(g^{-1})$ for any $g$ and $h$ in $G$.

The \textbf{degree of the representation} is the same as the dimension of the vector space.

\subsection{Tensor operations on representations}
	
	To-do here:
	\begin{enumerate}
		\setlength\itemsep{-1em}
		\item direct sums of representations
		\item definition of degree of representation
		\item definition of subrepresentation
	\end{enumerate}
	
\subsection{Irreducibility, Schur's and Maschke's}
	
	\begin{theorem}[Maschke's theorem]\label{thm:maschke}
		Let $G$ be a finite group and let $V$ be any arbitrary representation of $G$, then $V$ is composed of a number of not necessarily distinct irreducible subrepresentations $W_i$ of $V$, or in other words,
		\begin{align}
			V = \bigoplus_i W_i.
		\end{align}
	\end{theorem}
	
	\begin{note}
		The direct sum can be expressed in terms of \underline{distinct} subrepresentations $W_i$ of $V$ and their multiplicities $a_i$ as
		\begin{align}
			V = \bigoplus_i a_i W_i,
		\end{align}
		where $a_i W_i$ is the direct sum of $a_i$ copies of $W_i$.
	\end{note}
	
	To prove Maschke's theorem, first we establish a few definitions and lemmas.
	
	\begin{corollary}
		If the representation is of degree 1, $V = W_1$ and $\dim V = 1$, it is clearly irreducible.
	\end{corollary}
	
	\begin{definition}[Orthogonal subspace]
		\cite[12.7]{holst} If $W$ is a subspace in $V$, then its complementary subspace in $V$ is defined by 
		\[
			W^\perp = \left\lbrace \vvec \in V \ \middle\vert \ \wvec^\dagger \vvec = 0, \text{ for any } \wvec \in W \right\rbrace.
		\]
	\end{definition}
	
	\begin{lemma}[Complementary subspaces]\label{lemma:compsubspace}\cite[Thm 12.16]{holst}
		Let $W$ be a subspace of a space $V$. Let $\dim V = n$ and $\dim W = k$, then:
		\noindent
		\begin{enumerate}[i)]
			\setlength\itemsep{-0.5em}
			\item $\dim W^\perp = n - k$,
			\item $V = W \oplus W^\perp$,
			\item If $\{\ebas_i\}_{i=1}^{n}$ is the basis in $V$, then $\{\ebas_i\}_{i=1}^{k}$ and $\{\ebas_i\}_{i=k+1}^{n}$ are the bases for $W$ respectively $W^\perp$.
		\end{enumerate}
	\end{lemma}
	
	\begin{proof}[Proof of Maschke's theorem \ref{thm:maschke}]
		The theorem will be proved using induction on complementary subspaces. If $V$ itself is irreducible, we are done. If $V$ is not irreducible, then there is a non-trivial proper subrepresentation $W$ of $V$ with a basis $\{\ebas_i\}_i$, then according to lemma \ref{lemma:compsubspace}, there is a complementary ``orthogonal'' subrepresentation $W^\perp$ such that $V = W \oplus W^\perp$. If both of $W$ and $W^\perp$ are irreducible, then we are done. If either are not, then after recursion on $W$ and $W^\perp$ and since $V$ is a finite vector space, we are done.
	\end{proof}

\subsection{Some naïve examples}

In this section, $G$ is a finite group of order $|G| = n$. 

\paragraph*{The trivial and signature representations.}

	For any $G$, and choosing $V = \R$, a map
	\[
		\rho: G \rightarrow \GL(\R)
	\]
	defined by $ \rho(g) = 1$ for every $g$ in $G$, in other words mapping every element in $G$ to the trivial element of $V = \R$. It is trivially a homomorphism since $\rho(e)=1$ and both $\rho(gh)=1$ and $\rho(g)\rho(h)=1$ and is called the \textbf{trivial representation}. It is a representation of degree 1.
	
	%Letting $\mathbf{X}(g) = 1$ for every $g$ in $G$, we have a trivial representation of degree 1. This is trivially a homomorphism and thus a representation.
	
	Any group has a trivial degree 1 representation, but there are many trivial representations of larger degree, for example choosing $V = \R^3$, we define the representation by $\rho(g) = \1_3$, the three-by-three identity matrix. However, this is not an irreducible representation as $\R^3 = \R \oplus \R \oplus \R = 3\R$ and
	\[
		\begin{pmatrix} 1&0&0\\0&1&0\\0&0&1 \end{pmatrix} = (1)\oplus(1)\oplus(1).
	\]	
	
	The symmetric group has another degree 1 representation. Every element $\sigma$ of the symmetric group of degree $n$, denoted $S_n$, has a sign, or parity, denoted $\sgn(\sigma)$ defined by the number of transpositions\footnote{2-cycles} required to construct it. Letting $k$ be that number, then the parity or sign of $\sigma$ is $\sgn(\sigma)=(-1)^k$, so an \textbf{even} permutation has sign $+1$ and an \textbf{odd} permutation has sign $-1$. Now, with $G = S_n$ and choosing $V=\F_2$, this $\sgn$ function
	\[
		\sgn: S_n \rightarrow \{-1, 1\}
	\]
	defined as above is also a homomorphism since $\rho(e) = 1$ and for two arbitrary elements $\sigma$ and $\tau$ of $S_n$ with respective sign $(-1)^s$ and $(-1)^t$ we have 
	\[
		\rho(\sigma)\rho(\tau) = (-1)^s(-1)^t = (-1)^{st} = \rho(\sigma\tau).
	\] 
	This $\sgn$ function is called the \textbf{alternating representation}\footnote{In some literature called the sign or signature representation.} of the symmetric group.
	
	\paragraph{Permutation and regular representation}
	
	Let the group $G$ act on an appropriately chosen set $A$ by
	\[
		\phi : G \times A \rightarrow A
	\] % could be replaced by by ``permutating its elements''
	defined by $\phi: g \mapsto ga$. After constructing a vector space\marginnote{What is the dim?} $V$ with a basis $\mathcal{B} = ( \ebas_a )_{a \in A}$, we can define a representation
	\[
		\rho: G \rightarrow \GL(V)
	\]
	defined by
	\[
		\rho(g)(\ebas_a) = \ebas_{ga}.
	\]
	This is a homomorphism since for any $a$ in $A$, $\rho(e)$ acts as a identity map by 
	\[
		\rho(e)(\ebas_a) = \ebas_{ea} = \ebas_a
	\]
	and also for any $g$ and $h$ in $G$, 
	\[
		\rho(g)\rho(h)(\ebas_a) = \rho(g)(\ebas_{ha}) = \ebas_{gha} = \rho(gh)(\ebas_a).
	\]
	This homomorphism is called the \textbf{permutation representation}.
	
	\begin{example}[Some examples]
		Here are some suggestions for permutation representations
		\begin{itemize}
			\item $S_n$ can act on $\{1,2,\ \dots, n\}$, then $V$ is $\R^n$ and for a $\sigma$ in $S_n$, $\rho(\sigma)$ permutes the $i$th coordinate axis tho the $\sigma(i)$th axis.
			\item $Z_n$ (the cyclic group of size $n$) can act on the $n$th roots of unity (finite subset of $S_1$) by rotation???
			\item $\D_{2n}$ (the dihedral group of degree $n$, size $2n$) can act on the vertices of a regular $n$-gon (ostensibly the set $\{1, 2, \ \dots, n\}$)(I SUPPOSE??? HOW???)
		\end{itemize}
	\end{example}
	
	Instead of any general set, we can let $G$ act on itself. Similarly as above, we construct a vector space with basis $( \ebas_g )_{g \in G}$ and define a representation
	\[
		\rho: G \rightarrow \GL(V)
	\]
	defined by
	\[
		\rho(g)(\ebas_h) = \ebas_{gh}.
	\]
	To show that it is a homomorphism, one could follow the same steps as above. This is called the \textbf{regular representation}.
	
%\paragraph{The regular representation}
%
%In some sense, we can let $G$ act on itself by multiplication. Let $G$ be a finite group of order $n$ and $V$ be a vector space with basis $\left( \ebas_g \right)_{g \in G}$ defined by\footnote{Bad notation! $\ebas_1$ is the basis vector associated with the identity element $e$ of $G$.} $\ebas_1 = e$ and $\ebas_g = \rho(g)(\ebas_1)$. 
%
%\paragraph{The permutation representation}
%
%We can let $G$ act on any arbitrary set $A$, then,
%\[
%	\rho: G \times A \rightarrow A 
%\]
%defined by 
%\[
%	\rho(g): a \mapsto ga
%\]
%for any $g$ in $G$ and $a$ in $A$. Then $V$ is a vector space with basis $\left( \ebas_a \right)_{a \in A}$ where $\rho(g)(\ebas_a) = \ebas_{ga}$.
%
%It is ``natural'' to let $S_n$ act on $\{1, 2, \dots, n\}$.
%
%\begin{example}[Representations of $S_n$.]
%	Letting $\mathbf{T}(\sigma) = 1$ for every $\sigma$ in $S_n$, we have a trivial representation $T$ of $S_n$.
%	
%	Since $S_n$ is a group of permutations where every element has a sign property, the sign function
%	\[
%		\sgn(\sigma) = (-1)^k
%	\]
%	where $k$ is the least number of disjoint transpositions needed to construct $\sigma$, is an appropriate representation since $\sgn(\sigma)\sgn(\tau) = \sgn(\sigma\tau)$ for every $\sigma$ and $\tau$ in $S_n$.
%	
%	
%	
%	\begin{table}[hbt!]
%		\centering
%		\begin{tabular}{c | c c c c c c }
%			$S_3$  & $e$ & $(12)$ & $(13)$ & $(23)$ & $(123)$ & $(132)$ \\ \hline
%			 $T$   & 1   & 1      & 1      & 1      & 1       & 1       \\
%			$\sgn$ & 1   & -1     & -1     & -1     & 1       & 1       \\
%			 $R$   &     &        &        &        &         &         \\
%			 $P$   &     &        &        &        &         &
%		\end{tabular}
%	\end{table}
%	
%	
%\end{example} % RENSAD
%	\clearpage{\thispagestyle{empty}}
\section{Representation theory}

	\subsection{Introduction and basic definitions}
	
	A representation of a group in a vector space is a homomorphism that takes an element of the group into the group of automorphisms of the vector space. 
	
	For a vector space $V$ over a field $\K$, we denote its group of automorphisms $\GL(V)$. If $V$ is finite with $\dim V = n$ and a basis, usually denoted $\{\bas_i\}_{i=1}^n$, then $\GL(V)$ is identified with the set of linear transformations of $V$, which is the set of invertible square matrices of size $n$, denoted $\GL_n(\K)$\cite[18.1]{DummitFoote}. 
	\begin{definition}[Representation]
		Given a finite group $G$, a representation of $G$ in $V$ is a homomorphism
		\begin{align}
			\rho: G \rightarrow \GL(V).
		\end{align}
	\end{definition}
	In other words, $\rho$ takes an element of $G$ and maps it to a transformation of $\GL(V)$, or again $\rho(g): V \rightarrow V$ for some $g$ in $G$.	Recall that a homomorphism is defined by, for any two elements $g$ and $h$ in $G$, as
	\begin{align}\label{eq:homomorphism}
		\rho(gh) = \rho(g)\rho(h),
	\end{align}
	and two consequences follows:	
	\begin{proposition}
		For the identity element $e$ of $G$ and an arbitrary element $g$ with inverse $g^{-1}$ in $G$,
		\begin{align}
			\rho(e) = \id
		\end{align}
		and
		\begin{align}
			\rho(g)^{-1} = \rho\left(g^{-1}\right),
		\end{align}
		where $\id$ is the identity transformation and $\rho^{-1}(g)$ is the inverse representation of $g$.
	\end{proposition}
	\begin{proof}
		Take $g$ as an arbitrary element of $G$. The first identity follows from taking $h=e$ in eq.\ref{eq:homomorphism}:
		\[
			\rho(g) \overset{!}{=} \rho(eg) = \rho(e)\rho(g)
		\]
		and instead taking $h=g^{-1}$:
		\[
			\rho(g) \rho(g^{-1}) =  \rho(gg^{-1}) \overset{!}{=} \rho(e) = \id.
		\]
		we force the second identity.
	\end{proof}
	
	If the vector space $V$ is provided with a basis, then a matrix representations of $G$ may be found.
	\begin{definition}[Matrix representation]
		Given a finite group $G$, a matrix representation of $G$ in $V$ is a homomorphism
		\begin{align}
			X(g) \in \GL(V): V \rightarrow V.
		\end{align}
	\end{definition}
	
	\begin{note}
		Sometimes the vector space $V$ itself is said to be the representation, instead of the set of maps $\{\rho(g)\}_{g \in G}$.
	\end{note}
	
	\begin{definition}
		The degree of a representation is said to be the dimension of its vector space.
	\end{definition}
	
	\subsection{The trivial representation}
	
	Any group can be represented by a trivial representation of degree 1 defined by 
	\begin{align}
		\rho(g) = 1, \text{ for every } g \in G.
	\end{align}
	It is clearly a homomorphism\footnote{Here I did not state the vector space, should I? ``A vector space over the trivial group/field	''.}. 
	
	\subsection{The alternating representation of the symmetric group}
	
	Choosing $G$ to be the symmetric group of degree $n$, denoted $S_n$, we can find another degree 1 representation by studying the sign, or parity, of the elements of $S_n$. For any two permutations $\sigma$ and $\tau$ of $S_n$ with respective sign $(-1)^s$ and $(-1)^t$, the sign of their composition is 
	\begin{align}
		\sgn(\sigma \circ \tau) = (-1)^s \cdot (-1)^t = (-1)^{st} = \sgn(\sigma) \cdot \sgn(\tau),
	\end{align}
	so clearly this sign function is a homomorphism from $S_n$ to a one-dimensional vector space over the two elements\footnote{The vector space is over the field of two elements, $\F_2$.} $+1$ and $-1$, hence the $\sgn$ function is representation of degree 1 and is called the alternating representation, or alternatively the sign or signature representation.
	
	\subsection{The permutation representation}
	
	We could let the group $G$ act on an appropriately chosen set $X$, and study its action $f$ by permutation of the elements of $X$. This action is defined by any $g$ in $G$ and $x$ in $X$ as
	\begin{table}[hbt!]\centering\begin{tabular}{c c c c}
			$f:$ & $G \times X$    & $\rightarrow$ & $X$ \\
			&\rotatebox[origin=c]{90}{$\in$}&&\rotatebox[origin=c]{90}{$\in$} \\
			$f:$ & $(g,x)$ & $\mapsto$     & $gx$.
	\end{tabular}\end{table}

	Further more, define a vector space $V$ spanned by the orthogonal basis $\mathcal{B} = \{\ehat_x\}_x$ for every $x$ in $X$. Then we have a homomorphism
	\begin{align}
		\rho: G \rightarrow \GL(V)
	\end{align}
	defined by its action on the basis $\mathcal{B}$ by, for any $g$ in $G$,
	\begin{align}
		\rho(g): \ehat_x \mapsto g\ehat_x = \ehat_{gx}.
	\end{align}
	or on an arbitrary vector $v$ in $V$ as
	\begin{align}
		\rho(g): \vvec \mapsto g\vvec.
	\end{align}
	
	\textit{Show that it is a homomorphism!}
	
	\subsubsection{$S_n$ acting on $\N_n$}
		The action is defined as $S_n \times \N_n$ by $(\sigma,i) \mapsto \sigma(i)$. Choosing a basis $\mathcal{B} = \{\ehat_i\}_{i=1}^n$ we span a vector space $V$ of dimension $\dim V = n$. We then define a homomorphism $\rho$ from $S_n$ to $\GL(V)$ by
		\begin{align}
			\rho(\sigma): \ehat_i \mapsto \ehat_{\sigma(i)},
		\end{align}
		that is by permutating a basis vector $\ehat_i$ to another one $\ehat_j = \ehat_{\sigma(i)}$.
		
		\textit{action of $\rho(\sigma)$ on any vector??}
		
		The corresponding matrix representations are the permutation matrices
		\begin{align}
			X(\sigma) = (x_{ij})_{n \times n}, \text{ where } x_{ij} = \delta_{j,\sigma(i)} = \begin{cases}
				1 \text{ if } j = \sigma(i),\\
				0 \text{ otherwise.}
			\end{cases}
		\end{align}
		
		\begin{example}[$S_n$ for $n = 2,3$]
			
			The symmetric group of degree 2 has two elements, $S_2 = \{(1), (12)\}$ and their matrix representations in $\mathbf{\ehat_1},\mathbf{\ehat_2}$-space are
			\begin{table}[hbt!]
				\centering
				\caption{Matrix representations of $S_2$}
				\begin{tabular}{c | c c}
					$\sigma$    & $(1)$ & $(12)$ \\
					\hline
					$X(\sigma)$ & $\begin{pmatrix}
						1 & 0 \\ 0 & 1
					\end{pmatrix}$ & $\begin{pmatrix}
					0 & 1 \\ 1 & 0
					\end{pmatrix}$
				\end{tabular}
			\end{table}
			
			Likewise, representations of $S_3$ follows:
			
			\begin{table}[hbt!]
				\centering
				\caption{Matrix representations of $S_3$}
				\begin{tabular}{c | c c c c c c}
					$\sigma$  &$(1)$& $(12)$&$(13)$&$(23)$&$(123)$&$(132)$ \\
					\hline
					$X(\sigma)$ & 
					$\left(\begin{smallmatrix}
						1 & 0 & 0 \\
						0 & 1 & 0 \\
						0 & 0 & 1
					\end{smallmatrix}\right)$ &
					$\left(\begin{smallmatrix}
						0 & 1 & 0 \\
						1 & 0 & 0 \\
						0 & 0 & 1
					\end{smallmatrix}\right)$ &
					$\left(\begin{smallmatrix}
						0 & 0 & 1 \\
						0 & 1 & 0 \\
						1 & 0 & 0
					\end{smallmatrix}\right)$ &
					$\left(\begin{smallmatrix}
						1 & 0 & 0 \\
						0 & 0 & 1 \\
						0 & 1 & 0
					\end{smallmatrix}\right)$ &
					$\left(\begin{smallmatrix}
						0 & 0 & 1 \\
						1 & 0 & 0 \\
						0 & 1 & 0
					\end{smallmatrix}\right)$ &
					$\left(\begin{smallmatrix}
						0 & 1 & 0 \\
						0 & 0 & 1 \\
						1 & 0 & 0
					\end{smallmatrix}\right)$
				\end{tabular}
			\end{table}
			
%			$$S_3 = \{(1), (12), (13), (23), (123), (132)\}$$
%			$$S_4 = \{(1), 
%						(12), (13), (14), (23), (24), (34), 
%						(12)(34), (13)(24), (14)(23), 
%						(123), (124), (134), (132), (142), (143), (234), (243),
%						(1234), (1243), (1324), (1342), (1423), (1432)
%						\}$$
		\end{example}
		
	\subsection{The regular representation}
	
	Instead of letting the group $G$ act on some arbitrary set by permutation, we could let it act on itself. This action is defined as
	\[
		G \times G \rightarrow G
	\]
	which for every $g$ in $G$ is given by 
	\[
		(g,h) \mapsto gh,
	\]
	for a $h$ in $G$. The corresponding vector space $V$ is spanned by basis vectors $\{\ghat\}_{g \in G}$ constructed from the members of $G$. A representation of $G$ in $V$ is then
	\[
		\rho: G \rightarrow \GL(V)
	\]
	defined by
	\[
		\rho(g)(\hhat) = g\hhat.
	\] % RENSAD
	
	\clearpage{\thispagestyle{empty}}
\section{Character Theory}

%	For a representation $\rho$ of a group $G$ in a vector space $V$, the character of the representation is defined as the trace of the matrix representation in $\GL(V)$ of a element $g$ of $G$. If the matrix representation is given by
%	\begin{align}
%		X(g) = \left( x_{ij} \right)_{n \times n}
%	\end{align}
%	then the character is
%	\begin{align}
%		\chi_V(g) = \Tr X(g) = \sum_{i=1}^{n} x_{ii}.
%	\end{align}
%	Evident from this definition, the character of a degree 1 representation is the same as the representation, since a $1 \times 1$-matrix is usually identified with its only element.
%	
%	The character of the entire group $G$ in a representation $V$ is a vector 
%	\begin{align}
%		\chi_V = \left( \chi_V(g) \right)_{g \in G},
%	\end{align}
%	sometimes abbreviated to just containing a character for every conjugacy class
%	\begin{align}
%		\chi_V = \left( \chi_V([g]) \right)_{[g] \subseteq G},
%	\end{align}
%	since ``character'' is a class function on the elements of $G$.
%	
%	Since the trace of a matrix is invariant under conjugation, the character
	
	\begin{definition}[Character of a representation]\label{def:character}
		Let $\rho$ be a representation of degree $n$ of a group $G$ in a vector space $V$. The \textbf{character} of the representation is defined to be the trace of the representation. In other words, if $\rho(g)$ is the representation of a $g$ in $G$ and it has the matrix representation
		\begin{align}
			\rho_V(g) = (x_{ij})_{n \times n},
		\end{align}
		then the character of this $g$ is
		\begin{align}
			\chi_V(g) := \Tr \rho_V(g) = \sum_{i=1}^{n} x_{ii}.\label{eq:deftrace}
		\end{align}
		When $V$ is contextually obvious, it is omitted.
	\end{definition}
	
	\begin{corollary}[Character of degree 1 representation]\label{cor:chardeg1}
		Evident from definition \ref{def:character}, the character of a degree 1 representation is the same as the representation, since a $1 \times 1$-matrix is usually identified with its only element.
	\end{corollary}
	\begin{proof}
		For a degree 1 representation, $n=1$ and then equation \ref{eq:deftrace} becomes $\chi_V(g) = x_{11}$ which is identified with $(x_{11}) = \rho_V(g)$.
	\end{proof}
	
	Recall from linear algebra that two similar matrices have the same trace \cite[Thm 5.4.]{Holst} and that two matrices $X$ and $Y$ are similar if there exists another matrix $S$ such that $SXS^{-1} = Y$. 
	
	Also, recall that the trace is the sum of the eigenvalues of a matrix.
	
	\begin{lemma}[Trace function is a class function]
		The trace function = ``character function'' is invariant under conjugation, it is a class function on the elements of $G$, it partitions $G$ into disjoint conjugacy classes where all $g$ in the conjugacy class $[g]$ has the same character.
	\end{lemma}
	
	\begin{theorem}
		The character of a representation is independent of chosen matrix representation and chosen representative of a conjugacy class.
	\end{theorem}
	\begin{proof}
		Consequence of previous definitions and recollections.
	\end{proof}
	
	\begin{theorem}[Expected number of irreducible representations]\label{thm:numberirrep}
		The number of irreducible representations of a group is the same as the number of conjugacy classes of the group.
	\end{theorem}
	
	\subsection{Character tables of $Z_n$}

In section \ref{sect:reprZn}, the $n$ degree 1 representations of any cyclic group $Z_n$ was described, and by theorem \ref{thm:numberirrep}, those described are all irreducible representations of $Z_n$. Therefore the tables \ref{table:Z3}, \ref{table:Z4} and \ref{table:Z5} are the character tables of respectively $Z_3$, $Z_4$ and $Z_5$.

%In every character table, the first row corresponds to the trivial representation found in section \ref{sect:trivrepr} and for an even $n$, the character table also contains a ``alternating representation'' since these cyclic groups are isomorphic to the rotation subgroup of $\D_{2n}$.
	\subsection{Character table of $S_3$}

\begin{theorem}[Character of the permutation representation]
	The character of a permutation representation is the number of fixed points of the associated element of the symmetric group.
\end{theorem}
\begin{proof}
	The character is the sum of the diagonal elements of the matrix representation. The matrices of the permutation representations are permutation matrices, which have a 1 on the diagonal if that row represents a fixed point and 0 otherwise. Therefore the number of fixed point is the sum along the diagonal.
\end{proof}

The representations of $S_3$ found so far in sections \ref{sect:trivrepr} and \ref{sect:regS3} and examples \ref{ex:permS3} and \ref{ex:standS3} are presented in table \ref{table:charS3}. The group is presented in their conjugacy classes $[\sigma] \subset S_3$ where $\sigma$ is any representative of that class. Beneath every class is the number of elements in that class. Rows below the dashed horizontal line are not irreducible.

\begin{table}[hbt!]
	\centering
	\begin{tabular}{c | c c c}
		     $S_3$      & $[1]$   & $[12]$  & $[123]$ \\
		$\Ss|[\sigma]|$ & $\Ss 1$ & $\Ss 3$ & $\Ss 2$ \\ \hline
		   $\chi_T$     & 1       & 1       & 1       \\
		   $\chi_A$     & 1       & -1      & 1       \\
		   $\chi_S$     & 2       & 0       & -1      \\ \hline\hline
		   $\chi_P$     & 3       & 1       & 0       \\
		   $\chi_R$     & 6       & 0       & 0
	\end{tabular}
	\caption{Character table of $S_3$.}
	\label{table:charS3}
\end{table}

\paragraph{Decomposition of $P$.} In section $\ref{sect:standSn}$, the permutation representation was decomposed through geometric arguments to the direct sum of the trivial and the standard representation, this is seen in the table as $\chi_T + \chi_S = (1,1,1) + (2,0,-1) = (3,1,0) = \chi_P$, confirming that $P = T \oplus S$. 

\paragraph{Decomposition of $R$.} In section \ref{sect:regS3}, the matrices of the regular representation was exhaustively calculated. The character of $R$ thus is $\chi_R = (6,0,0)$, since the identity element is the only element in $S_3$ that leaves every other element fixed. In fact this is true for any $S_n$:
\begin{theorem}[Character of the regular representation of $S_n$]\label{thm:charregSn}
	 The character of the regular representation for any $S_n$ is $(n!, 0, \dots, 0)$, where $n!$ is the size of $S_n$.
\end{theorem}
\begin{proof}
	The center of $S_n$ is trivial.
\end{proof}

To decompose $R$, we take the inner product of $\chi_R$ with every other irreducible character.
\begin{align}
	\bra{\chi_T}\ket{\chi_R} = 1, \quad
	\bra{\chi_A}\ket{\chi_R} = 1, \quad
	\bra{\chi_S}\ket{\chi_R} = 2, 
\end{align}
so $\chi_R = \chi_T + \chi_A + 2\chi_S$ and $R = T \oplus A \oplus 2S$.

Also:
\begin{theorem}
	If $V_i$ are the irreducible representations of a group, then the regular representation $R$ is the direct sum of $\dim V_i$ copies of $V_i$, ie. 
	\begin{align}
		R = \bigoplus_{i=1}^{k} (\dim V_i)V_i = \bigoplus_{i=1}^{k} \chi_{V_i}(e)V_i .
	\end{align}
\end{theorem}
\begin{proof}
	Needs a proof and/or source, noticed while staring at character tables.
\end{proof}

\paragraph{Decomposition of $S \otimes S$.} The character of $S \otimes S$ is $\chi_{S \otimes S} = \chi_S \cdot \chi_S = (4,0,1)$, which by a quick glance on the character table is seen to be the sum of all irreducibles, ie. $\chi_{S \otimes S} = \chi_T + \chi_A + \chi_S$ and $S \otimes S = T \oplus A \oplus S$.

\paragraph{Decomposition of $S^{\otimes n}$.}\cite[Exercise 2.7.]{FultonHarris} To find the decomposition of larger tensor powers than $n=2$, we study the character $\chi_{S^{\otimes n}} = (2^n, 0, (-1)^n)$ and take the inner product of it with the other irreducibles and find:
\begin{align}
	\bra{\chi_T}\ket{\chi_{S^{\otimes n}}} = \bra{\chi_A}\ket{\chi_{S^{\otimes n}}} = \frac{1}{6}\left(2^n + (-1)^n\right), \\
	\bra{\chi_S}\ket{\chi_{S^{\otimes n}}} = \frac{1}{6}\left(2^{n+1} + (-1)^{n+1}\right), 
\end{align}
so $\chi_{S^{\otimes n}} = a_n \chi_T \oplus a_n \chi_A + a_{n+1} \chi_S$ and thus $S^{\otimes n} = a_n T \oplus a_n A \oplus a_{n+1} S$ where $a_n = \left(2^n + (-1)^n\right)/6$.

\paragraph{Decomposition of $\Sym^2S$ and $\bigwedge^2S$.} The character of $\Sym^2S$ is $\chi_{\Sym^2S} = \frac{1}{2}\left[ (4,0,1) - (2,2,-1) \right] = (1,-1,1) = \chi_A$, so $\Sym^2S \cong A$. The character of $\bigwedge^2S$ is $\chi_{\bigwedge^2S} = (3,1,0) = \chi_P = \chi_T + \chi_S$, so $\bigwedge^2S \cong P \cong T \oplus S$, confirming $S \otimes S = \Sym^2S \oplus \bigwedge^2S$ since $\chi_{\Sym^2S} + \chi_{\bigwedge^2S} = A \oplus T \oplus S = \chi_{S \otimes S}$. Also, $R = S \otimes S \oplus S$.

\paragraph{} The contributions from the last few paragraphs are to the now complete character table of $S_3$ is presented in table \ref{table:completecharS3}.

\begin{table}[hbt!]
	\centering
	
	\begin{tabular}{c | c c c | l}
		       $S_3$         & $[1]$   & $[12]$  & $[123]$ &                                          \\
		  $\Ss|[\sigma]|$    & $\Ss 1$ & $\Ss 3$ & $\Ss 2$ &  \textit{Alternate compositions}         \\ \hline
		      $\chi_T$       & 1       & 1       & 1       &                                          \\
		      $\chi_A$       & 1       & -1      & 1       & $ A \cong \Sym^2S$                       \\
		      $\chi_S$       & 2       & 0       & -1      &                                          \\ \hline\hline
		      $\chi_P$       & 3       & 1       & 0       & $ P \cong T \oplus S \cong \bigwedge^2S$ \\
		$\chi_{S \otimes S}$ & 4       & 0       & 1       & $ S \otimes S \cong T \oplus A \oplus S$ \\
		      $\chi_R$       & 6       & 0       & 0       & $ R \cong T \oplus A \oplus 2S$
	\end{tabular}
	\label{table:completecharS3}
	\caption{Complete character table of $S_3$. The representations above the dashed line are irreducibles, and those below are composed. Some compositions are presented in the right-most column.}
\end{table}


	\subsection{Character table of $S_4$}

This section follows the same methods from the previous section on $S_3$. So far, we have three irreducible representations of $S_4$: The trivial and the alternating representation from section \ref{sect:trivrepr}; the standard representation from section \ref{sect:standSn} and also we can surmise that there is a regular representation of $S_4$ with character $(24,0,0,0,0)$ from theorem \ref{thm:charregSn}. Along with the composite permutation representation, the character table of $S_4$ is presented in table \ref{table:charS4}. 

\begin{table}[hbt!]
	\centering
	\begin{tabular}{c | c c c c c}
		     $S_3$      & $[1]$   & $[12]$  & $[123]$ & $[1234]$ & $[(12)(34)]$ \\
		$\Ss|[\sigma]|$ & $\Ss 1$ & $\Ss 6$ & $\Ss 8$ & $\Ss 6$  & $\Ss 3$      \\ \hline
		   $\chi_T$     & 1       & 1       & 1       & 1        & 1            \\
		   $\chi_A$     & 1       & -1      & 1       & -1       & 1            \\
		   $\chi_S$     & 3       & 1       & 0       & -1       & -1           \\ \hdashline
		   $\chi_P$     & 4       & 2       & 1       & 0        & 0            \\
		   $\chi_R$     & 24      & 0       & 0       & 0        & 0            	\end{tabular}
	\caption{Character table of $S_4$.}
	\label{table:charS4}
\end{table}

Recall that there are five conjugacy classes in $S_4$, then by theorem \ref{thm:numberirrep} we can expect two more irreducible representations of $S_3$. To find them, we compose new representations out of the ones we have found so far, using the toolbox of section \ref{sect:tensorops}.

\paragraph{Decomposition of $A \otimes S$.} The character of this composition is $(3,-1,0,1,-1)$, an ``alternating'' version of $S$, and taking the inner product of it with itself we find $|\chi_{A \otimes S}|^2 = 1$ and thus $A \otimes S$ is an irreducible representation. From this point on it is denoted as $S'$.

\paragraph{Decomposition of $S \otimes S$.} The character of this composition is $(9,1,0,1,1)$ and it is found to be linearly dependent on $\chi_T$,  $\chi_S$ and $\chi_{A \otimes S}$ with overlaps of 1 and independent of $\chi_A$, however $T$, $S$ and $S'$ together are of degrees 1, 3 and 3 and $S \otimes S$ is of degree 9, so there is either another representation of degree 2 or two of degree 1. \textit{SOURCE, theorem om square sum of $\dim V_i$}. Denoting this representation by $V$, its character is $\chi_V = \chi_{S \otimes S} - \chi_T - \chi_S - \chi_{S'} = (2, 0, -1, 0, 2)$, which is found to be such that $|\chi_V|^2 = 1$, hence $V$ is irreducible and $S \otimes S$ is found to be decomposed to $T \oplus V \oplus S \oplus S'$. We have now found all five irreducible representations of $S_4$.

\paragraph{Decomposition of $R$.} The character of $R$ is $(24,0,0,0,0)$ and its overlaps with $T, A, V, S$ and $S'$ are respectively 1, 1, 2, 3 and 3, so $R = T \oplus A \oplus 2V \oplus 3S \oplus 3S'$.

\paragraph{Decomposition of $\Sym^2S$ and $\bigwedge^2S$.} The character of $\Sym^2S$ is 
\begin{align}
	\frac{1}{2}\left[ (9,1,0,1,1)-(3,3,0,-1,3) \right] = (3,-1,0,1,-1)
\end{align}
which is the character of $S'$. Likewise, the character of $\bigwedge^2S$ is $(6,2,0,0,2)$ which is found to be the sum of the characters of $T, S$ and $V$, hence $\bigwedge^2S = T \oplus S \oplus V$.

\paragraph{Decomposition of the $n$th tensor power of $V$.}

The character of $V^{\otimes n}$ is
\begin{align}
	\chi_{V^{\otimes n}} = (2^n, 0, (-1)^n, 0, 2^n).
\end{align}

Its overlap with $T, A, V, S$ and $S'$ are respectively $a_n, a_n, a_{n+1}, 0$ and $0$, where 
\begin{align}
	a_n := \frac{1}{24}\left( 4 \cdot 2^n + 8 \cdot (-1)^n \right),
\end{align}
so the decomposition is
\begin{align}
	V^{\otimes n} = a_n (T \oplus A) \oplus a_{n+1} V.
\end{align}

\paragraph{Decomposition of the $n$th tensor power of $S$.} The character of $S^{\otimes n}$ is 
\begin{align}
	(\chi_S)^n = (3^n, 1, 0, (-1)^n, (-1)^n),
\end{align}
and after projecting it on the characters of the irreducible representations from table \ref{table:completecharS4}, we find that
\begin{align}\label{eq:nthpowerS}
	S^{\otimes n} = a_n T \oplus b_n A \oplus c_n V \oplus a_{n+1} S \oplus b_{n+1} S',
\end{align}
where 
\begin{align}
	\begin{cases}
		a_n = \frac{1}{24}(3^n + 9(-1)^n+6), \\
		b_n = \frac{1}{24}(3^n - 3(-1)^n -6), \text{ and} \\
		c_n = \frac{1}{24}(2\cdot 3^n + 6(-1)^n).
	\end{cases}
\end{align}

\paragraph{Decomposition of the $n$th tensor powers of $S'$.} Since $S' = A \otimes S$, the $n$th power of $S'$ is the tensor product of $n$ factors of $A \otimes S$. Since the tensor product is a commutative operation, we can rearrange it as
\begin{align}
	(A \otimes S)^{\otimes n} = S^{\otimes n} \otimes A^{\otimes n}
\end{align}
The $n$th tensor power of $A$ is $T$ if $n$ is even and $A$ if $n$ is odd, so for even $n$, $S'^{\otimes n} \cong S^{\otimes n}$, and for odd $n$, $S'^{\otimes n} \cong S^{\otimes n} \otimes A$. Then the character is 
\begin{align}
	\chi_{S'^{\otimes n}} = \begin{cases}
		\big(3^n,1,0,(-1)^{n}, (-1)^n\big), \text{ if } n \text{ is even, or} \\
		\big(3^n,-1,0,(-1)^{n+1}, (-1)^n\big), \text{ if } n \text{ is odd.}
	\end{cases}
\end{align}
Tensor multiplying equation \ref{eq:nthpowerS} with $A$, we have:
\begin{align}
	S^{\otimes n} \otimes A &= \Big(a_n T \oplus b_n A \oplus c_n V \oplus a_{n+1} S \oplus b_{n+1} S'\Big)\otimes A \\
	&= b_n T \oplus a_n A \oplus c_n V \oplus b_{n+1} S \oplus a_{n+1} S',
\end{align}
and then we have found the decomposition of all tensor powers of $S'$. For even $n$, it is the same as $S^{\otimes n}$, for odd $n$, $A$ and $T$, and $S$ and $S'$ switch multiplicities.

\paragraph{} The findings of the last few paragraphs (except those on larger tensor powers) are presented in table \ref{table:completecharS4}.

\begin{table}[hbt!]
	\centering
	\begin{tabular}{c | c c c c c | l}
		       $S_4$         & $[1]$   & $[12]$  & $[123]$ & $[1234]$ & $[(12)(34)]$ &                                                              \\
		  $\Ss|[\sigma]|$    & $\Ss 1$ & $\Ss 6$ & $\Ss 8$ & $\Ss 6$  & $\Ss 3$      & \textit{Alternate compositions}                              \\ \hline
		      $\chi_T$       & 1       & 1       & 1       & 1        & 1            &                                                              \\
		      $\chi_A$       & 1       & -1      & 1       & -1       & 1            &                                                              \\
		      $\chi_V$       & 2       & 0       & -1      & 0        & 2            &                                                              \\
		      $\chi_S$       & 3       & 1       & 0       & -1       & -1           &                                                              \\
		$\chi_{S'}$ & 3       & -1      & 0       & 1        & -1           & $S' :\cong A \otimes S \cong \Sym^2S$                                  \\ \hline\hline
		      $\chi_P$       & 4       & 2       & 1       & 0        & 0            & $P \cong T \oplus S$                                         \\
		   $\chi_{\bigwedge^2S}$    & 6       & 2       & 0       & 0        & 2            & $\bigwedge^2S \cong T \oplus S \oplus V$                     \\
		$\chi_{S \otimes S}$ & 9       & 1       & 0       & 1        & 1            & $S \otimes S \cong T \oplus V \oplus S \oplus S' \cong S' \otimes S'$ \\
		      $\chi_R$       & 24      & 0       & 0       & 0        & 0            & $R \cong T \oplus A \oplus 2V \oplus 3S \oplus 3S'$
	\end{tabular}
	\caption{Complete character table of $S_4$. The representations above the dashed line are irreducibles, and those below are composed. Some compositions are presented in the right-most column.}
	\label{table:completecharS4}
\end{table}
	\subsection{Character table of the Klein 4-group $V$}

The Klein 4-group, denoted by\footnote{$V$ is for vierer, four in german.} $V$ is, along with $Z_4$, the only group of order 4. It can be described as the set
\begin{align}
	V = \{ e, x, y, xy \},
\end{align}
where every element is its own inverse and the product of any two distinct non-identity elements is the third one.
It is an abelian group, which means that every element has its own conjugacy class (SOURCE).

Since every group has a trivial representation $T$ in any vector space, so do $V$. The center of $V$ is trivial, so the regular representation $R$ has the character $(4,0,0,0)$ and is clearly not irreducible.

If $V$ instead is considered as a subgroup of $S_4$, that is
\begin{align}
	V = \{ (1), (12)(34), (13)(24), (14)(23) \} < S_4,
\end{align}
then each entry in the character of the permutation representation $P$ is the number of fixed points of the associated element. Then one sees that the character of the permutation representation is also $(4,0,0,0)$, it isomorphic to the regular representation.

From theorem \ref{thm:numberirrep} we expect three additional irreducible representations. Since the square sum over all representations of the character of $e$ must be equal to the order of the group (SOURCE), they must be of degree 1. We know that they must be parallel to themselves and orthogonal to $T$ and have a 1 as the character of $e$, so the three additional are alterations of $T$ with two of the ones exchanged with negative ones. Then $R$ is the direct sum of these four irreducibles. These findings are presented in table \ref{table:charV} as the representations $A, B$ and $C$.


\begin{table}[hbt!]
	\centering
	\begin{tabular}{c | c c c c l}
		$V$ & $e$ & $x$ & $y$ & $xy$ &             \\ \hline
		$T$ & 1   & 1   & 1   & 1   &             \\
		$A$ & 1   & -1  & -1  & 1   &             \\
		$B$ & 1   & -1  & 1   & -1  &             \\
		$C$ & 1   & 1   & -1  & -1  &             \\ \hline\hline
		$R$ & 4   & 0   & 0   & 0   & %$R \cong P \cong T \oplus A \oplus B \oplus C$
	\end{tabular}
	\caption{Character table for $V$.}
	\label{table:charV}
\end{table}
	\subsection{Character table of the dihedral group of order 8, $\D_8$}

$\D_8$ is a 2-Sylow subgroup of $S_4$. Restrict the character table of $S_4$ to $\D_8$. 

$\D_8$ has its non-trivial                                                                                                                                                                                                                                                              center $Z(\D_8) = \{e, r^2\}$ as a normal subgroup. The quotient $\D_8/\{e, r^2\}$ is isomorphic to the Klein 4-group. The character table of the Klein 4-group is a ``subset'' of the character table of $\D_8$.


Hej Erik $a+b=c$
	
%	\clearpage{\thispagestyle{empty}}
\section{Conclusion}

A conclusion will be written here.  \lipsum[8]

Further topics. \lipsum[9]

Alternate paths (modules) \lipsum[10] % Ej skriven


	

	%%The section with references.
%
\clearpage{\thispagestyle{empty}}
%\bibliographystyle{alphaurl} %This style prints the urls of online sources.
\bibliographystyle{alpha}
%
%
%% All references are to be stored in the file bibliographyFile.bib! 
%% It is in the same folder as your main .tex file.
%% It is using the bibtex syntax, and you can get these entries from google scholar, etc.
%% Your supervisor should know,
%
\bibliography{bibliography} 
%

\end{document}
